\documentclass[11pt]{article}
\usepackage[margin=1in]{geometry}
\usepackage[T1]{fontenc}
\usepackage[utf8]{inputenc}
\usepackage[french]{babel}
\usepackage{amsmath,amssymb,amsthm}
\usepackage[colorlinks=true,linkcolor=blue,urlcolor=blue]{hyperref}
\newtheorem{problem}{Problème}
\newenvironment{refs}{\par\small\noindent\emph{Références :}\begin{itemize}\setlength\itemsep{0pt}}{\end{itemize}\normalsize}
\title{Hors de la Caverne \\ \Large Analyse réelle et théorie de la mesure}
\author{}
\date{}
\begin{document}
\maketitle
\noindent\emph{Des problèmes, pas de réponses.} Chaque problème ci-dessous est posé
sans solution. Les références indiquent où trouver les connaissances nécessaires.
\medskip


\section*{Lycée}

\begin{problem}[{Bornes supérieures --- Lycée}]
\textbf{RA-01.} \textbf{Problème.} Pour un ensemble $ S $ de nombres réels, un \textit{majorant}
est un nombre $ b $ tel que $ x \le b $ pour tout $ x \in S $ ; la \textit{borne supérieure}
$ \sup S $ est le plus petit des majorants, et la \textit{borne inférieure} $ \inf S $ le plus
grand des minorants. Déterminez, avec démonstration, la borne supérieure et la borne inférieure de
chacun des ensembles suivants, et décidez dans chaque cas si elles appartiennent à l'ensemble :
\begin{enumerate}
\item[(a)] $ \{ 1/n : n = 1, 2, 3, \dots \} $ ;
\item[(b)] $ \{ n/(n+1) : n = 1, 2, 3, \dots \} $ ;
\item[(c)] $ \{ x \in \mathbb{Q} : x > 0 \text{ et } x^2 < 2 \} $.
Pour (c), expliquez pourquoi la borne supérieure existe dans $ \mathbb{R} $ alors que l'ensemble
n'a pas de borne supérieure à l'intérieur de $ \mathbb{Q} $.
\end{enumerate}

\end{problem}

\noindent L'\textit{axiome de la borne supérieure} --- toute partie non vide et majorée de $\mathbb{R}$
admet un plus petit majorant --- est l'unique propriété qui sépare $\mathbb{R}$ de $\mathbb{Q}$, et toute l'analyse repose
sur elle. Richard Dedekind l'a isolée en 1858 en préparant un cours de calcul différentiel à
Zurich, mécontent du flou géométrique entretenu autour de la droite numérique, et il a publié en
1872 sa construction de $\mathbb{R}$ par \og{} coupures \fg{}. La question (c) est exactement le trou de $\mathbb{Q}$ que
$ \sqrt{2} $ vient combler.

\begin{refs}
\item Contexte : Borne supérieure et borne inférieure --- Wikipédia (\url{https://fr.wikipedia.org/wiki/Borne_sup%C3%A9rieure_et_borne_inf%C3%A9rieure})
\item Contexte : Complétude des nombres réels --- Wikipédia (en anglais) (\url{https://en.wikipedia.org/wiki/Completeness_of_the_real_numbers})
\item Lecture : Abbott, \textit{Understanding Analysis}, ch. 1
\end{refs}

\bigskip

\begin{problem}[{Pourquoi 0,999\dots{} = 1 --- Lycée}]
\textbf{RA-02.} \textbf{Problème.} Le symbole $ 0.999\ldots $ désigne, par définition, la
limite de la suite $ 0.9,\ 0.99,\ 0.999,\ \dots $ --- autrement dit la somme de la série
$ \frac{9}{10} + \frac{9}{100} + \frac{9}{1000} + \cdots $. Démontrez que
$ 0.999\ldots = 1 $. Donnez deux démonstrations : l'une en sommant la série géométrique, l'autre
directement à partir de la définition d'une limite (montrez que pour tout $ \varepsilon > 0 $,
les termes de la suite finissent par se trouver à une distance inférieure à $ \varepsilon $ de 1).
Expliquez ensuite précisément ce qui cloche dans l'objection \og{} $ 0.999\ldots $ est plus petit que
1 parce qu'il n'atteint jamais 1 \fg{}.
\end{problem}

\noindent C'est l'égalité la plus discutée d'internet, et la discussion porte toujours sur
les définitions : un développement décimal illimité \textit{est} une limite, et une fois cela admis
la démonstration est courte. La question est réellement instructive --- elle force la prise de
conscience, rendue précise pour la première fois par Cauchy dans son \textit{Cours d'analyse} (1821),
que les processus infinis ne prennent sens qu'à travers la notion de limite. C'est aussi une
première rencontre avec le fait qu'un nombre réel peut avoir deux écritures décimales.

\begin{refs}
\item Contexte : 0,999\dots{} --- Wikipédia (\url{https://fr.wikipedia.org/wiki/0,999%E2%80%A6})
\item Lecture : Abbott, \textit{Understanding Analysis}, ch. 2
\end{refs}

\bigskip

\begin{problem}[{La série harmonique diverge --- Lycée}]
\textbf{RA-03.} \textbf{Problème.} Démontrez que la série harmonique
$ 1 + \frac{1}{2} + \frac{1}{3} + \frac{1}{4} + \cdots $ diverge : les sommes partielles
$ H_n = 1 + \frac{1}{2} + \cdots + \frac{1}{n} $ dépassent tout majorant. Montrez ensuite
quantitativement que $ H_{2^k} \ge 1 + k/2 $, et servez-vous-en pour estimer combien de termes
sont nécessaires avant que les sommes partielles ne dépassent 10 pour la première fois.
\end{problem}

\noindent Nicole Oresme a trouvé l'argument classique de regroupement vers 1350 --- une
démonstration médiévale que l'on ne peut toujours pas améliorer en élégance ---, après quoi elle fut
oubliée puis redémontrée par Pietro Mengoli (1650) et Jakob Bernoulli (1689). Le résultat est un
avertissement permanent : les termes d'une série peuvent tendre vers zéro pendant que la somme
croît sans limite. La question compagne, la valeur exacte de $ \sum 1/n^2 $, est le problème de
Bâle, traité dans L'art de la démonstration (\url{proofs.html}).

\begin{refs}
\item Contexte : Série harmonique --- Wikipédia (\url{https://fr.wikipedia.org/wiki/S%C3%A9rie_harmonique})
\item Lecture : Abbott, \textit{Understanding Analysis}, ch. 2
\item Cours : MIT OCW 18.100A Real Analysis (\url{https://ocw.mit.edu/courses/18-100a-real-analysis-fall-2020/})
\end{refs}

\bigskip

\begin{problem}[{Les deux limites qu'une suite possède toujours --- Lycée, short}]
\textbf{RA-24.} \textbf{Problème.} Pour la suite $ a_n = (-1)^n \left( 1 + \frac{1}{n}
\right) $, déterminez $ \limsup_{n \to \infty} a_n $ et $ \liminf_{n \to \infty} a_n $, avec
démonstration. Démontrez ensuite qu'une suite bornée de nombres réels converge si et seulement si sa
limite supérieure et sa limite inférieure coïncident.
\end{problem}

\noindent Une suite bornée n'a aucune raison de converger, mais elle possède toujours une
plus grande et une plus petite \og{} valeur finale \fg{} --- et, contrairement à la limite, ces deux-là
existent toujours. Cauchy travaillait déjà avec la plus grande valeur d'adhérence d'une suite ; en
faire une définition plutôt qu'une façon de parler est ce qui permet à l'analyse d'énoncer des
théorèmes (la règle de Cauchy, le lemme de Fatou) valables sans aucune hypothèse de convergence.

\begin{refs}
\item Contexte : Limite supérieure et limite inférieure --- Wikipédia (\url{https://fr.wikipedia.org/wiki/Limite_sup%C3%A9rieure_et_limite_inf%C3%A9rieure})
\item Lecture : Abbott, \textit{Understanding Analysis}, ch. 2
\end{refs}

\bigskip

\begin{problem}[{D'où vient e --- Lycée, short}]
\textbf{RA-25.} \textbf{Problème.} Démontrez que la suite
$ a_n = \left( 1 + \frac{1}{n} \right)^{n} $ est strictement croissante et majorée par
$ 3 $, de sorte qu'elle converge ; sa limite est par définition le nombre $ e $. Déduisez-en
que $ 2 < e < 3 $.
\end{problem}

\noindent Jacob Bernoulli a rencontré cette suite en 1683 en se demandant ce qu'il advient
des intérêts composés lorsqu'on les compose de plus en plus souvent ; les intérêts croissent, mais
pas sans limite, et le plafond est $ e $. Euler a donné à la constante sa lettre vers 1731 et l'a
calculée avec dix-huit décimales. La démonstration est une pure application du principe selon lequel
une suite croissante et majorée converge --- le premier endroit où un nombre réel est \textit{créé} par
une limite au lieu d'être écrit.

\begin{refs}
\item Contexte : e (nombre) --- Wikipédia (\url{https://fr.wikipedia.org/wiki/E_(nombre)})
\item Lecture : Rudin, \textit{Principles of Mathematical Analysis}, ch. 3
\end{refs}

\bigskip

\begin{problem}[{Segments emboîtés, et comment Cantor a dénombré la droite --- Lycée}]
\textbf{RA-26.} \textbf{Problème.} Ne supposez de $ \mathbb{R} $ rien d'autre que d'être un corps totalement
ordonné dans lequel toute partie non vide majorée admet un plus petit majorant (RA-01).
\begin{enumerate}
\item[(a)] Démontrez le théorème des segments emboîtés : si $ I_1 \supseteq I_2 \supseteq I_3 \supseteq \cdots $
sont des intervalles fermés bornés $ I_n = [a_n, b_n] $, alors $ \bigcap_{n=1}^{\infty} I_n $ n'est
pas vide. Montrez par un exemple que la conclusion peut tomber en défaut si les intervalles sont
ouverts, ou s'ils ne sont pas bornés.
\item[(b)] Déduisez-en le théorème de Bolzano--Weierstrass : toute suite bornée de nombres réels admet une
sous-suite convergente.
\item[(c)] Démontrez que $ \mathbb{R} $ n'est pas dénombrable, en utilisant (a) plutôt que les
développements décimaux : étant donné une suite quelconque $ x_1, x_2, x_3, \dots $ de nombres
réels, construisez des intervalles fermés emboîtés $ I_1 \supseteq I_2 \supseteq \cdots $ avec
$ x_n \notin I_n $, et concluez qu'aucune suite ne peut énumérer tous les nombres réels.
\item[(d)] Démontrez que les nombres algébriques --- les racines des polynômes non nuls à coefficients
entiers --- forment un ensemble dénombrable. Combinez avec (c) pour en déduire que les nombres
transcendants existent. Remarquez bien que cet argument n'en nomme aucun.
\end{enumerate}

\end{problem}

\noindent L'article de Georg Cantor de 1874, \og{} Sur une propriété de la collection de tous
les nombres algébriques réels \fg{}, contenait exactement l'argument des questions (c) et (d) --- le
célèbre argument diagonal n'est venu que dix-sept ans plus tard, en 1891. Trente ans auparavant,
Liouville avait construit explicitement un nombre transcendant, à la main, au prix d'un effort
énorme ; Cantor a montré en deux pages que presque tout nombre est transcendant sans en exhiber un
seul. Kronecker trouvait ce raisonnement répugnant et tenta d'en bloquer la publication. La question
(a) est l'humble outil qui fait tout fonctionner, et c'est aussi la forme la plus limpide de la
complétude : c'est elle qui distingue $ \mathbb{R} $ de $ \mathbb{Q} $.

\begin{refs}
\item Contexte : Théorème des fermés emboîtés --- Wikipédia (\url{https://fr.wikipedia.org/wiki/Th%C3%A9or%C3%A8me_des_ferm%C3%A9s_embo%C3%AEt%C3%A9s})
\item Contexte : Argument de la diagonale de Cantor --- Wikipédia (\url{https://fr.wikipedia.org/wiki/Argument_de_la_diagonale_de_Cantor})
\item Contexte : Nombre transcendant --- Wikipédia (\url{https://fr.wikipedia.org/wiki/Nombre_transcendant})
\item Lecture : Abbott, \textit{Understanding Analysis}, ch. 1
\end{refs}

\bigskip

\begin{problem}[{La propriété d'Archimède, et comment $\mathbb{Q}$ remplit $\mathbb{R}$ --- Lycée}]
\textbf{RA-36.} \textbf{Problème.} Ne supposez de $ \mathbb{R} $ rien d'autre que d'être un corps totalement
ordonné dans lequel toute partie non vide majorée admet un plus petit majorant.
\begin{enumerate}
\item[(a)] Démontrez la propriété d'Archimède : $ \mathbb{N} $ n'est pas majoré dans $ \mathbb{R} $,
donc pour tout réel $ x $ il existe un entier positif $ n > x $. (Supposez $ \mathbb{N} $
majoré et examinez $ \sup \mathbb{N} - 1 $.) Déduisez-en que pour tout
$ \varepsilon > 0 $ il existe un $ n $ tel que $ 1/n < \varepsilon $.
\item[(b)] Démontrez que les rationnels sont denses : dès que $ a < b $ il existe un rationnel $ q $
tel que $ a < q < b $. (Choisissez $ n $ tel que $ 1/n < b-a $ et prenez le plus petit
entier $ m $ tel que $ m/n > a $ ; expliquez pourquoi un tel plus petit entier existe.)
\item[(c)] Démontrez que $ \sqrt{2} $ est irrationnel, et déduisez-en que les irrationnels sont denses
eux aussi.
\item[(d)] Démontrez que tout réel $ x > 0 $ possède exactement une racine $ n $-ième positive, en
montrant que $ y = \sup\{ t > 0 : t^n < x \} $ vérifie $ y^n = x $ --- éliminez tour à tour
$ y^n < x $ et $ y^n > x $. Signalez chaque endroit où (a) est utilisé.
\item[(e)] Montrez que (a) ne découle pas des seuls axiomes de corps totalement ordonné : ordonnez le corps
des fractions rationnelles $ \mathbb{R}(t) $ en déclarant $ p/q > 0 $ lorsque les coefficients
dominants de $ p $ et de $ q $ sont de même signe, vérifiez que cela en fait un corps totalement
ordonné, et constatez que $ t $ y est plus grand que tout entier.
\end{enumerate}

\end{problem}

\noindent La propriété de (a) est plus ancienne que son nom : Archimède l'énonce comme un
postulat dans \textit{De la sphère et du cylindre} et en attribue l'idée à Eudoxe, qui s'en servit pour
rendre rigoureuse la méthode d'exhaustion au quatrième siècle avant notre ère. La question (e) est la
raison pour laquelle il s'agit d'un véritable théorème sur $ \mathbb{R} $ et non d'une trivialité ---
il existe de parfaits corps totalement ordonnés contenant des éléments infiniment grands, et c'est la
complétude qui les exclut ici. Tout ce qui, en analyse, commence par \og{} choisissons $ n $ assez
grand \fg{} est une application de (a), et (d) est l'endroit où le nombre $ \sqrt{2} $ que RA-01 avait
localisé comme borne supérieure devient enfin un nombre.

\begin{refs}
\item Contexte : Archimédien (propriété d'Archimède) --- Wikipédia (\url{https://fr.wikipedia.org/wiki/Archim%C3%A9dien})
\item Contexte : Partie dense --- Wikipédia (\url{https://fr.wikipedia.org/wiki/Partie_dense})
\item Contexte : Racine d'un nombre --- Wikipédia (\url{https://fr.wikipedia.org/wiki/Racine_d'un_nombre})
\item Lecture : Rudin, \textit{Principles of Mathematical Analysis}, ch. 1
\item Lecture : Abbott, \textit{Understanding Analysis}, ch. 1
\end{refs}

\bigskip

\begin{problem}[{Une récurrence qui monte vers une limite --- Lycée, short}]
\textbf{RA-37.} \textbf{Problème.} Soient $ a_1 = \sqrt{2} $ et $ a_{n+1} = \sqrt{2 + a_n} $
pour $ n \ge 1 $. Démontrez que la suite est strictement croissante et majorée par $ 2 $, donc
convergente, et déterminez sa limite --- en justifiant soigneusement pourquoi la limite doit vérifier
l'équation que vous écrivez, et pourquoi cette équation n'a qu'une seule racine admissible.
\end{problem}

\noindent La démonstration, c'est le théorème de convergence monotone et rien d'autre, mais
c'est dans la dernière clause que se trouve la réflexion : passer à la limite à l'intérieur de
$ \sqrt{2+\cdot} $ utilise la continuité, et écarter la seconde racine de l'équation du second
degré obtenue utilise l'encadrement que vous avez démontré. Ces radicaux imbriqués sont aussi le
mécanisme de la formule de Viète pour $\pi$ (1593), le premier produit infini jamais écrit, où les mêmes
expressions apparaissent comme cosinus d'angles divisés en deux encore et encore.

\begin{refs}
\item Contexte : Théorème de convergence monotone --- Wikipédia (\url{https://fr.wikipedia.org/wiki/Th%C3%A9or%C3%A8me_de_convergence_monotone})
\item Contexte : Radical imbriqué --- Wikipédia (\url{https://fr.wikipedia.org/wiki/Radical_imbriqu%C3%A9})
\item Lecture : Abbott, \textit{Understanding Analysis}, ch. 2
\end{refs}

\bigskip


\section*{Licence}

\begin{problem}[{$\varepsilon$--$\delta$, en pratique --- Licence}]
\textbf{RA-04.} \textbf{Problème.} Une fonction $ f $ est continue en un point $ c $ si
pour tout $ \varepsilon > 0 $ il existe un $ \delta > 0 $ tel que
$ |x - c| < \delta $ entraîne $ |f(x) - f(c)| < \varepsilon $.
\begin{enumerate}
\item[(a)] Démontrez directement à partir de cette définition que $ f(x) = x^2 $ est continue en tout
$ c \in \mathbb{R} $.
\item[(b)] Démontrez que $ g(x) = 1/x $ est continue en tout point de
$ (0, \infty) $, mais qu'elle n'y est \textit{pas uniformément} continue --- c'est-à-dire qu'aucun
$ \delta $ unique ne peut convenir à tous les points à la fois pour un $ \varepsilon $ donné.
\item[(c)] Démontrez que
$ g $ \textit{est} uniformément continue sur $ [1, \infty) $.
\end{enumerate}

\end{problem}

\noindent Bolzano (1817) et Cauchy (1821) ont formulé la continuité au moyen d'inégalités,
et les cours berlinois de Weierstrass des années 1860 ont fixé le rituel $\varepsilon$--$\delta$ que tout analyste
exécute aujourd'hui. La distinction faite en (b) et (c) --- continuité en chaque point contre
continuité uniforme sur un ensemble --- a mis des décennies à être vue clairement, et la confusion
entre les deux a invalidé plusieurs \og{} théorèmes \fg{} publiés au début du dix-neuvième siècle.

\begin{refs}
\item Contexte : Fonction continue --- Wikipédia (\url{https://fr.wikipedia.org/wiki/Fonction_continue})
\item Contexte : Continuité uniforme --- Wikipédia (\url{https://fr.wikipedia.org/wiki/Continuit%C3%A9_uniforme})
\item Lecture : Abbott, \textit{Understanding Analysis}, ch. 4
\item Cours : MIT OCW 18.100A Real Analysis (\url{https://ocw.mit.edu/courses/18-100a-real-analysis-fall-2020/})
\end{refs}

\bigskip

\begin{problem}[{Deux fonctions monstrueuses --- Licence}]
\textbf{RA-05.} \textbf{Problème.}
\begin{enumerate}
\item[(a)] Soit $ D(x) = 1 $ si $ x $ est rationnel et
$ D(x) = 0 $ si $ x $ est irrationnel (la fonction de Dirichlet). Démontrez que $ D $ n'est
continue en aucun point de $ \mathbb{R} $.
\item[(b)] Soit $ T(x) = 1/q $ lorsque $ x = p/q $ est
rationnel écrit sous forme irréductible avec $ q > 0 $, et $ T(x) = 0 $ lorsque $ x $ est
irrationnel (la fonction de Thomae). Démontrez que $ T $ est continue en tout point irrationnel et
discontinue en tout point rationnel.
\end{enumerate}

\end{problem}

\noindent Dirichlet a produit $ D $ en 1829, dans son article sur les séries de Fourier,
comme premier exemple d'une fonction défiant toute intuition fondée sur les formules ; la variante de
Thomae est apparue en 1875. Ensemble, elles marquent le moment où \og{} fonction \fg{} s'est mise à
signifier une règle arbitraire plutôt qu'une expression analytique. Savoir si l'image miroir de (b)
peut exister --- une fonction continue exactement sur les rationnels --- trouve sa réponse dans le
théorème de Baire, en RA-13.

\begin{refs}
\item Contexte : Fonction de Dirichlet --- Wikipédia (\url{https://fr.wikipedia.org/wiki/Fonction_de_Dirichlet})
\item Contexte : Fonction de Thomae --- Wikipédia (\url{https://fr.wikipedia.org/wiki/Fonction_de_Thomae})
\item Lecture : Abbott, \textit{Understanding Analysis}, ch. 4
\end{refs}

\bigskip

\begin{problem}[{Suites de Cauchy et complétude --- Licence}]
\textbf{RA-06.} \textbf{Problème.} Une suite $ (a_n) $ est \textit{de Cauchy} si pour tout
$ \varepsilon > 0 $ il existe un $ N $ tel que $ |a_m - a_n| < \varepsilon $ pour
tous $ m, n \ge N $.
\begin{enumerate}
\item[(a)] Démontrez que toute suite convergente est de Cauchy.
\item[(b)] En utilisant la propriété de la
borne supérieure (via le théorème de Bolzano--Weierstrass, que vous démontrerez également : toute
suite bornée admet une sous-suite convergente), démontrez la réciproque : toute suite de Cauchy de
nombres réels converge.
\item[(c)] Exhibez une suite de Cauchy de nombres rationnels sans limite
rationnelle, et concluez que $ \mathbb{Q} $ n'est pas complet.
\end{enumerate}

\end{problem}

\noindent Cauchy a énoncé le critère en 1821 mais n'a pas pu démontrer qu'il est suffisant ---
il n'existait encore aucune construction de $\mathbb{R}$ à partir de laquelle le démontrer. Elle est venue en
1872, quand Cantor (à la suite de Méray, 1869) a construit les nombres réels \textit{comme} classes
d'équivalence de suites de Cauchy de rationnels, tandis que Dedekind les construisait à partir de
coupures. Ce problème parcourt le cercle historique : la complétude est exactement ce qu'il faut
ajouter à $\mathbb{Q}$ pour que le critère fonctionne.

\begin{refs}
\item Contexte : Suite de Cauchy --- Wikipédia (\url{https://fr.wikipedia.org/wiki/Suite_de_Cauchy})
\item Contexte : Théorème de Bolzano--Weierstrass --- Wikipédia (\url{https://fr.wikipedia.org/wiki/Th%C3%A9or%C3%A8me_de_Bolzano-Weierstrass})
\item Lecture : Rudin, \textit{Principles of Mathematical Analysis}, ch. 3
\end{refs}

\bigskip

\begin{problem}[{Réarranger l'infini : le théorème de réarrangement de Riemann --- Licence}]
\textbf{RA-07.} \textbf{Problème.}
\begin{enumerate}
\item[(a)] Démontrez que la série harmonique alternée
$ 1 - \frac{1}{2} + \frac{1}{3} - \frac{1}{4} + \cdots $ converge, mais pas absolument.
\item[(b)] Démontrez le théorème de réarrangement de Riemann : si une série de nombres réels converge
conditionnellement (elle converge, mais pas absolument), alors pour tout nombre réel $ L $ il
existe un réarrangement de ses termes qui converge vers $ L $, et il existe des réarrangements
divergeant vers $ +\infty $ et vers $ -\infty $.
\item[(c)] Démontrez au contraire que tout réarrangement d'une série absolument convergente converge vers
la même somme.
\end{enumerate}

\end{problem}

\noindent Dirichlet a remarqué en 1837 que réarranger la série harmonique alternée peut en
changer la somme ; Riemann, dans son mémoire d'habilitation de 1854 sur les séries trigonométriques
(publié en 1867), a transformé l'observation en théorème définitif. C'est l'avertissement le plus
tranchant qui soit sur le fait que l'addition infinie n'est pas commutative, et la raison pour
laquelle la \og{} convergence absolue \fg{} mérite sa place centrale en analyse.

\begin{refs}
\item Contexte : Théorème de réarrangement de Riemann --- Wikipédia (\url{https://fr.wikipedia.org/wiki/Th%C3%A9or%C3%A8me_de_r%C3%A9arrangement_de_Riemann})
\item Lecture : Rudin, \textit{Principles of Mathematical Analysis}, ch. 3
\item Lecture : Abbott, \textit{Understanding Analysis}, ch. 2
\end{refs}

\bigskip

\begin{problem}[{Le théorème des valeurs intermédiaires --- Licence}]
\textbf{RA-08.} \textbf{Problème.} Démontrez le théorème des valeurs intermédiaires : si $ f $ est
continue sur $ [a, b] $ et si $ y $ est compris entre $ f(a) $ et $ f(b) $, alors
$ f(c) = y $ pour un certain $ c \in [a, b] $. Donnez une démonstration utilisant la borne
supérieure de $ \{ x : f(x) \le y \} $, et signalez où la complétude de $ \mathbb{R} $ est
utilisée. Déduisez-en ensuite :
\begin{enumerate}
\item[(a)] tout polynôme de degré impair à coefficients réels admet une racine réelle ;
\item[(b)] toute
fonction continue $ f : [0, 1] \to [0, 1] $ admet un point fixe ;
\item[(c)] à tout instant, il existe
deux points antipodaux de l'équateur terrestre où la température est la même (modélisez la
température par une fonction continue sur un cercle).
\end{enumerate}

\end{problem}

\noindent L'opuscule de Bolzano de 1817 --- \og{} une démonstration purement analytique du
théorème selon lequel entre deux valeurs donnant des résultats de signes opposés se trouve au moins
une racine réelle \fg{} --- est l'acte de naissance de l'analyse rigoureuse : il tenait à ce qu'un énoncé
sur les fonctions continues méritât une démonstration exempte de tout recours géométrique. Le
théorème est faux sur $\mathbb{Q}$, donc toute démonstration correcte doit invoquer la complétude ; observer
l'endroit où elle intervient est tout l'intérêt de l'exercice.

\begin{refs}
\item Contexte : Théorème des valeurs intermédiaires --- Wikipédia (\url{https://fr.wikipedia.org/wiki/Th%C3%A9or%C3%A8me_des_valeurs_interm%C3%A9diaires})
\item Lecture : Abbott, \textit{Understanding Analysis}, ch. 4
\item Cours : MIT OCW 18.100A Real Analysis (\url{https://ocw.mit.edu/courses/18-100a-real-analysis-fall-2020/})
\end{refs}

\bigskip

\begin{problem}[{Le théorème des bornes atteintes --- Licence}]
\textbf{RA-09.} \textbf{Problème.} Démontrez le théorème des bornes atteintes : une fonction
continue $ f $ sur un intervalle fermé borné $ [a, b] $ est bornée et atteint son maximum et son
minimum. (Démontrez d'abord qu'elle est bornée, par l'absurde avec Bolzano--Weierstrass ; puis qu'elle
atteint ses bornes.) Montrez par des exemples que chaque hypothèse est nécessaire : exhibez des
fonctions continues non bornées sur $ (0, 1) $ et sur $ [0, \infty) $, ainsi qu'une fonction
continue bornée sur $ (0, 1) $ n'atteignant ni maximum ni minimum.
\end{problem}

\noindent Weierstrass a démontré le théorème dans ses cours berlinois des années 1860, et
c'est la raison pour laquelle \og{} fermé et borné \fg{} --- la compacité, comme dirait le vingtième siècle ---
est devenu un concept porteur. L'existence de maxima avait été supposée en silence pendant deux
siècles, notamment dans des démonstrations fautives du principe de Dirichlet ; l'insistance de
Weierstrass à exiger ici une démonstration a remodelé le calcul des variations.

\begin{refs}
\item Contexte : Théorème des valeurs extrêmes --- Wikipédia (\url{https://fr.wikipedia.org/wiki/Th%C3%A9or%C3%A8me_des_valeurs_extr%C3%AAmes})
\item Lecture : Rudin, \textit{Principles of Mathematical Analysis}, ch. 4
\end{refs}

\bigskip

\begin{problem}[{La convergence simple n'est pas la convergence uniforme --- Licence}]
\textbf{RA-10.} \textbf{Problème.} Une suite de fonctions $ f_n $ converge
\textit{simplement} vers $ f $ si $ f_n(x) \to f(x) $ pour chaque $ x $ fixé, et
\textit{uniformément} si $ \sup_x |f_n(x) - f(x)| \to 0 $.
\begin{enumerate}
\item[(a)] Démontrez qu'une limite uniforme de fonctions continues est continue.
\item[(b)] Trouvez une suite de fonctions continues sur
$ [0, 1] $ convergeant simplement vers une fonction discontinue --- il existe un exemple classique
en une ligne ; démontrez que votre convergence n'est pas uniforme.
\item[(c)] Trouvez des fonctions continues $ f_n $ sur
$ [0, 1] $ telles que $ f_n(x) \to 0 $ pour tout $ x $ alors que $ \int_0^1 f_n = 1 $ pour
tout $ n $ : la convergence simple ne commute pas avec l'intégration.
\end{enumerate}

\end{problem}

\noindent Cauchy affirmait en 1821 qu'une série convergente de fonctions continues a une
somme continue --- faux tel quel, comme Abel le fit remarquer en 1826 avec une série de Fourier. La
réparation, la notion de convergence uniforme, a été distillée par Seidel, Stokes et surtout
Weierstrass dans les années 1840--1850. Les contre-exemples demandés ici sont ceux que tout analyste
transporte avec lui ; l'exemple (c) est la \og{} bosse qui s'échappe \fg{}, qui motivera plus tard le
théorème de convergence dominée (RA-16).

\begin{refs}
\item Contexte : Convergence uniforme --- Wikipédia (\url{https://fr.wikipedia.org/wiki/Convergence_uniforme})
\item Lecture : Abbott, \textit{Understanding Analysis}, ch. 6
\item Lecture : Rudin, \textit{Principles of Mathematical Analysis}, ch. 7
\end{refs}

\bigskip

\begin{problem}[{Le monstre de Weierstrass --- Licence}]
\textbf{RA-11.} \textbf{Problème.} (Problème de compréhension.) La fonction de Weierstrass est
\[ W(x) = \sum_{n=0}^{\infty} a^n \cos(b^n \pi x), \]
où $ 0 < a < 1 $ et où $ b $ est un entier impair vérifiant $ ab > 1 + 3\pi/2 $. Elle
est continue partout et dérivable nulle part.
\begin{enumerate}
\item[(a)] Démontrez vous-même la moitié facile : $ W $ est
continue sur $ \mathbb{R} $ (démontrez et utilisez le test M de Weierstrass).
\item[(b)] Lisez une démonstration de la non-dérivabilité partout et rédigez-en un compte rendu dans vos
propres mots : quel rôle joue la condition $ ab > 1 $, pourquoi les fréquences $ b^n $
doivent-elles croître si vite (lacunarité), et où l'argument produit-il deux points voisins dont les
taux d'accroissement divergent violemment ?
\end{enumerate}

\end{problem}

\noindent Weierstrass a présenté l'exemple à l'Académie de Berlin en 1872 (publié par du
Bois-Reymond en 1875), démolissant la croyance universelle selon laquelle une fonction continue doit
être dérivable sauf en des points isolés. Hermite écrivait qu'il se détournait \og{} avec effroi et
horreur de cette plaie lamentable des fonctions sans dérivées \fg{}. Le théorème de Baire (RA-13) a
livré plus tard l'ironie finale : de telles fonctions ne sont pas des monstres mais le cas
typique.

\begin{refs}
\item Contexte : Fonction de Weierstrass --- Wikipédia (\url{https://fr.wikipedia.org/wiki/Fonction_de_Weierstrass})
\item Lecture : Abbott, \textit{Understanding Analysis}, ch. 6
\item Lecture : Stein \& Shakarchi, \textit{Real Analysis} (Princeton Lectures in Analysis III)
\end{refs}

\bigskip

\begin{problem}[{L'ensemble de Cantor --- Licence}]
\textbf{RA-12.} \textbf{Problème.} Soit $ C $ l'ensemble triadique de Cantor : de
$ [0, 1] $ on retire le tiers médian ouvert $ (1/3, 2/3) $, puis le tiers médian de chaque
intervalle restant, et ainsi de suite ; $ C $ est ce qui subsiste après toutes les étapes.
Démontrez :
\begin{enumerate}
\item[(a)] $ C $ est fermé, ne contient aucun intervalle, et sa \og{} longueur \fg{} est nulle (la longueur
totale retirée vaut 1) ;
\item[(b)] $ C $ n'est pas dénombrable --- en fait $ C $ est exactement l'ensemble des nombres de
$ [0, 1] $ admettant un développement en base 3 n'utilisant que les chiffres 0 et 2, ce qui donne
une correspondance de type bijection avec toutes les suites binaires infinies ;
\item[(c)] tout point de $ C $
est limite d'autres points de $ C $ ($ C $ est parfait). Concluez qu'un ensemble peut être
négligeable en longueur tout en étant aussi grand que $ \mathbb{R} $ en cardinal.
\end{enumerate}

\end{problem}

\noindent La construction a été publiée par Henry J. S. Smith en 1874 et, avec bien plus
d'influence, par Cantor en 1883, au cours de ses recherches sur les ensembles exceptionnels des
séries trigonométriques. C'est la drosophile de l'analyse réelle : la source standard des
contre-exemples, la première fractale (de dimension $ \log 2 / \log 3 $) et --- via la fonction de
Cantor, l'\og{} escalier du diable \fg{} qui lui est associé --- une fonction continue qui monte de 0 à 1
pendant que sa dérivée s'annule presque partout.

\begin{refs}
\item Contexte : Ensemble de Cantor --- Wikipédia (\url{https://fr.wikipedia.org/wiki/Ensemble_de_Cantor})
\item Contexte : Escalier de Cantor --- Wikipédia (\url{https://fr.wikipedia.org/wiki/Escalier_de_Cantor})
\item Lecture : Abbott, \textit{Understanding Analysis}, ch. 3
\end{refs}

\bigskip

\begin{problem}[{Catégorie de Baire --- Licence}]
\textbf{RA-13.} \textbf{Problème.} Un ensemble est \textit{d'intérieur vide} (rare) si son adhérence ne contient
aucun intervalle ouvert ; un ensemble est \textit{maigre} (de première catégorie) s'il est réunion
dénombrable d'ensembles rares.
\begin{enumerate}
\item[(a)] Démontrez le théorème de Baire pour $ \mathbb{R} $ : la droite réelle n'est pas
maigre --- de façon équivalente, une intersection dénombrable d'ouverts denses de $ \mathbb{R} $ est
dense.
\item[(b)] Déduisez-en que $ \mathbb{R} $ n'est pas dénombrable.
\item[(c)] Déduisez-en qu'il n'existe aucune fonction
$ f : \mathbb{R} \to \mathbb{R} $ continue en tout point rationnel et discontinue en tout point
irrationnel (utilisez le fait que l'ensemble des points de continuité d'une fonction quelconque est
une intersection dénombrable d'ouverts).
\end{enumerate}

\end{problem}

\noindent René-Louis Baire a démontré le théorème dans sa thèse de doctorat de 1899,
créant une notion de \og{} petitesse \fg{} pour les ensembles qui est topologique et non métrique --- et
exaspérément indépendante de la mesure : $\mathbb{R}$ se scinde en un ensemble maigre et un ensemble
négligeable. Les arguments de catégorie sont devenus l'un des grands moteurs non constructifs de
l'analyse, prouvant l'existence d'objets (comme les fonctions nulle part dérivables, qui forment un
ensemble co-maigre de $ C[0,1] $) sans jamais en exhiber un seul.

\begin{refs}
\item Contexte : Théorème de Baire --- Wikipédia (\url{https://fr.wikipedia.org/wiki/Th%C3%A9or%C3%A8me_de_Baire})
\item Contexte : Ensemble maigre --- Wikipédia (\url{https://fr.wikipedia.org/wiki/Ensemble_maigre})
\item Lecture : Oxtoby, \textit{Measure and Category}
\end{refs}

\bigskip

\begin{problem}[{Le théorème de Dini --- Licence, short}]
\textbf{RA-27.} \textbf{Problème.} Soit $ K $ un espace métrique compact et soient
$ f_n : K \to \mathbb{R} $ continues avec $ f_n(x) \le f_{n+1}(x) $ pour tout $ n $ et tout
$ x $, convergeant simplement vers une fonction continue $ f $. Démontrez que la convergence est
en fait uniforme, et montrez par un contre-exemple que la conclusion tombe en défaut si l'on
supprime soit la compacité de $ K $, soit la continuité de la limite.
\end{problem}

\noindent RA-10 montre que la convergence simple est bien plus faible que la convergence
uniforme ; le théorème de Dini identifie une paire d'hypothèses supplémentaires peu coûteuses ---
monotonie et compacité --- qui comble l'écart, et la démonstration est une unique et élégante
application du sous-recouvrement fini. Ulisse Dini l'a publié dans ses \textit{Fondamenti per la teorica
delle funzioni di variabili reali} de 1878, l'un des premiers livres entièrement écrits dans le
nouveau style rigoureux, et c'est encore la manière standard de renforcer une convergence dans les
arguments d'approximation.

\begin{refs}
\item Contexte : Théorèmes de Dini --- Wikipédia (\url{https://fr.wikipedia.org/wiki/Th%C3%A9or%C3%A8mes_de_Dini})
\item Lecture : Rudin, \textit{Principles of Mathematical Analysis}, ch. 7
\end{refs}

\bigskip

\begin{problem}[{La convexité impose la continuité --- Licence, short}]
\textbf{RA-28.} \textbf{Problème.} Dites que $ f : (a,b) \to \mathbb{R} $ est \textit{convexe} si
$ f(t x + (1-t) y) \le t f(x) + (1-t) f(y) $ pour tous $ x, y \in (a,b) $ et tout
$ t \in [0,1] $. Démontrez qu'une fonction convexe sur un intervalle ouvert est continue --- et même
lipschitzienne sur tout sous-intervalle fermé --- et montrez par un exemple qu'une fonction convexe
sur un intervalle fermé $ [a,b] $ peut être discontinue en une extrémité.
\end{problem}

\noindent C'est l'un des rares endroits de l'analyse où une inégalité purement algébrique
impose une conclusion topologique, sans qu'aucune continuité ne soit supposée nulle part. Tout
découle de l'inégalité des trois cordes --- les pentes des cordes croissent quand on se déplace vers
la droite ---, qui donne aussi des dérivées à droite et à gauche en tout point et la dérivabilité en
dehors d'un ensemble dénombrable. Johan Jensen a fait de la convexité un sujet à part entière en
1906 ; l'échec aux extrémités est précisément la raison pour laquelle l'analyse convexe moderne
travaille sur des domaines ouverts, ou avec des fonctions semi-continues inférieurement à valeurs
dans la droite réelle achevée.

\begin{refs}
\item Contexte : Fonction convexe --- Wikipédia (\url{https://fr.wikipedia.org/wiki/Fonction_convexe})
\item Lecture : Rudin, \textit{Real and Complex Analysis}, ch. 3
\item Lecture : R. T. Rockafellar, \textit{Convex Analysis}, partie I
\end{refs}

\bigskip

\begin{problem}[{Variation bornée et décomposition de Jordan --- Licence}]
\textbf{RA-29.} \textbf{Problème.} Pour $ f : [a,b] \to \mathbb{R} $, la \textit{variation totale} est
\[ V_a^b(f) \;=\; \sup \sum_{i=1}^{n} \left| f(x_i) - f(x_{i-1}) \right|, \]
la borne supérieure étant prise sur toutes les subdivisions $ a = x_0 < x_1 < \cdots < x_n = b $.
Dites que $ f $ est \textit{à variation bornée} si $ V_a^b(f) < \infty $.
\begin{enumerate}
\item[(a)] Démontrez que les fonctions à variation bornée sur $ [a,b] $ forment un espace vectoriel, que
$ V_a^b(f) = V_a^c(f) + V_c^b(f) $ dès que $ a < c < b $, et que toute fonction monotone
et toute fonction lipschitzienne est à variation bornée.
\item[(b)] Démontrez le théorème de décomposition de Jordan : $ f $ est à variation bornée si et
seulement si $ f = g - h $ pour deux fonctions croissantes au sens large $ g, h $ sur
$ [a,b] $.
\item[(c)] Déduisez-en qu'une fonction à variation bornée admet une limite à gauche et une limite à droite
en tout point, et n'a qu'un nombre au plus dénombrable de discontinuités.
\item[(d)] Montrez que $ f(x) = x \sin(1/x) $ pour $ x \in (0,1] $ avec $ f(0) = 0 $ est continue
mais pas à variation bornée, tandis que $ g(x) = x^2 \sin(1/x) $ avec $ g(0) = 0 $ est à
variation bornée. Expliquez ce que ce contraste dit sur les courbes continues qui ont une longueur
finie.
\end{enumerate}

\end{problem}

\noindent Camille Jordan a introduit la variation bornée en 1881 dans un but concret :
élargir la classe des fonctions dont on sait que la série de Fourier converge, en prolongeant le
théorème de Dirichlet au-delà des fonctions monotones par morceaux. La question (b) révèle ce que la
définition dit vraiment --- une fonction à variation bornée est une fonction monotone moins une autre,
si bien que toute la théorie de la dérivabilité des fonctions monotones lui est transférée d'un
coup. La variation totale est littéralement la distance verticale parcourue par un point mobile,
raison pour laquelle elle définit la longueur d'une courbe rectifiable ; et sa descendante en
théorie de la mesure --- les fonctions dont la dérivée est une mesure --- est aujourd'hui le cadre
standard du débruitage d'images et des ondes de choc des lois de conservation.

\begin{refs}
\item Contexte : Fonction à variation bornée --- Wikipédia (\url{https://fr.wikipedia.org/wiki/Fonction_%C3%A0_variation_born%C3%A9e})
\item Contexte : Variation totale d'une fonction --- Wikipédia (\url{https://fr.wikipedia.org/wiki/Variation_totale_d'une_fonction})
\item Lecture : Royden \& Fitzpatrick, \textit{Real Analysis}, 4e éd., ch. 6
\item Lecture : Wheeden \& Zygmund, \textit{Measure and Integral}, ch. 2
\end{refs}

\bigskip

\begin{problem}[{L'escalier du diable --- Licence}]
\textbf{RA-38.} \textbf{Problème.} Soit $ C \subset [0,1] $ l'ensemble triadique de Cantor de RA-12. Définissez
la fonction de Cantor $ F : [0,1] \to [0,1] $ comme suit : si $ x \in C $ a pour développement
triadique $ x = \sum_{n\ge1} 2c_n 3^{-n} $ avec tous les $ c_n \in \{0,1\} $, posez
$ F(x) = \sum_{n\ge1} c_n 2^{-n} $ ; sur chaque intervalle ouvert retiré lors de la construction
de $ C $, posez $ F $ constante, égale à sa valeur aux extrémités.
\begin{enumerate}
\item[(a)] Démontrez que $ F $ est bien définie (les deux développements triadiques d'une extrémité
donnent la même valeur), croissante au sens large et surjective sur $ [0,1] $ ; déduisez-en que
$ F $ est continue, et même uniformément continue, en utilisant seulement le fait qu'une
surjection croissante sur un intervalle ne peut avoir aucun saut.
\item[(b)] Démontrez que $ F'(x) = 0 $ en tout $ x $ hors de $ C $, donc que $ F' = 0 $ presque
partout, alors que $ F(1)-F(0) = 1 $. Concluez que $ F $ est continue, monotone et dérivable
presque partout, sans pour autant être l'intégrale de sa dérivée : le théorème fondamental de
l'analyse est mis en défaut pour elle de la façon la plus radicale possible.
\item[(c)] Démontrez que $ F $ n'est pas lipschitzienne, et trouvez son exposant de Hölder exact :
montrez que $ |F(x)-F(y)| \le c\,|x-y|^{\alpha} $ vaut avec $ \alpha = \log 2/\log 3 $ et une
certaine constante $ c $, et qu'elle échoue pour tout exposant plus grand.
\item[(d)] Posez $ G(x) = \frac{x + F(x)}{2} $. Démontrez que $ G $ est un homéomorphisme strictement
croissant de $ [0,1] $ sur lui-même avec $ m\big(G(C)\big) = \tfrac12 $. En prenant une partie
non mesurable $ A \subseteq G(C) $ du type construit en RA-18, démontrez que $ G^{-1}(A) $ est
un ensemble Lebesgue-mesurable (et même négligeable) dont l'image par l'homéomorphisme $ G $ n'est
pas mesurable --- et déduisez-en que $ G^{-1}(A) $ est Lebesgue-mesurable mais non borélien.
\end{enumerate}

\end{problem}

\noindent Cantor a introduit cette fonction en 1884, dans le même cercle de travaux que
l'ensemble de RA-12 ; Scheeffer puis Lebesgue l'ont étudiée, et les physiciens l'ont rebaptisée
escalier du diable après avoir rencontré la même forme dans le verrouillage de fréquences
d'oscillateurs couplés. C'est le contre-exemple qui impose toute la théorie de la continuité absolue
(RA-40) : une fonction peut s'élever de 0 à 1 tout en restant immobile presque partout, si bien que
\og{} l'intégrale de la dérivée redonne la fonction \fg{} requiert une hypothèse à laquelle personne n'avait
songé. La question (d) est la démonstration standard du fait que la tribu de Lebesgue est
strictement plus grande que la tribu borélienne, et elle montre que la mesurabilité n'est pas
préservée par un changement de variable continu.

\begin{refs}
\item Contexte : Escalier de Cantor --- Wikipédia (\url{https://fr.wikipedia.org/wiki/Escalier_de_Cantor})
\item Contexte : Fonction singulière --- Wikipédia (\url{https://fr.wikipedia.org/wiki/Fonction_singuli%C3%A8re})
\item Lecture : Stein \& Shakarchi, \textit{Real Analysis}, ch. 3
\item Lecture : Gelbaum \& Olmsted, \textit{Counterexamples in Analysis}
\end{refs}

\bigskip

\begin{problem}[{Une dérivée ne peut pas sauter --- Licence, short}]
\textbf{RA-39.} \textbf{Problème.} Démontrez le théorème de Darboux : si $ f $ est dérivable
en tout point de $ [a,b] $, alors $ f' $ prend toute valeur comprise entre $ f'(a) $ et
$ f'(b) $ --- bien que $ f' $ n'ait aucune raison d'être continue. Exhibez ensuite une fonction
dérivable dont la dérivée est discontinue en un point, et déduisez-en qu'une fonction présentant une
discontinuité de saut n'est jamais la dérivée de quoi que ce soit.
\end{problem}

\noindent Gaston Darboux a démontré cela dans son mémoire de 1875 sur les fonctions
discontinues, et c'est le premier signe que la classe des dérivées est une classe étrange : une
dérivée n'a pas à être continue, mais elle ne peut pas non plus être arbitraire. La question
naturelle qui suit --- quelles fonctions \textit{sont} des dérivées ? --- se révèle n'avoir aucune réponse
connue, et c'est le problème ouvert de RA-45.

\begin{refs}
\item Contexte : Théorème de Darboux (analyse) --- Wikipédia (\url{https://fr.wikipedia.org/wiki/Th%C3%A9or%C3%A8me_de_Darboux_(analyse)})
\item Lecture : Rudin, \textit{Principles of Mathematical Analysis}, ch. 5
\end{refs}

\bigskip


\section*{Master}

\begin{problem}[{Riemann contre Lebesgue --- Master}]
\textbf{RA-14.} \textbf{Problème.}
\begin{enumerate}
\item[(a)] Démontrez que la fonction de Dirichlet de RA-05 n'est pas
Riemann-intégrable sur $ [0, 1] $, tandis que la fonction de Thomae l'est, d'intégrale nulle.
\item[(b)] Démontrez
le critère de Lebesgue : une fonction bornée sur $ [a, b] $ est Riemann-intégrable si et seulement
si son ensemble de discontinuités est de mesure de Lebesgue nulle.
\item[(c)] Expliquez, en une page, le changement de
point de vue de l'intégrale de Lebesgue --- subdiviser l'image plutôt que le domaine --- et
vérifiez que la fonction de Dirichlet est Lebesgue-intégrable, d'intégrale nulle.
\end{enumerate}

\end{problem}

\noindent Riemann a défini son intégrale dans le mémoire de 1854 déjà rencontré en RA-07 et
s'est demandé au juste à quel point une fonction intégrable peut être discontinue ; la thèse de
Lebesgue de 1902, \og{} Intégrale, longueur, aire \fg{}, y a répondu et a remplacé l'intégrale purement et
simplement. L'image de Lebesgue lui-même pour (c) : un commerçant peut compter la recette du jour
pièce par pièce dans l'ordre où elles sont arrivées (Riemann), ou bien trier d'abord les pièces par
valeur faciale (Lebesgue) --- et la seconde méthode survit à des passages à la limite qui détruisent
la première.

\begin{refs}
\item Contexte : Intégrale de Riemann --- Wikipédia (\url{https://fr.wikipedia.org/wiki/Int%C3%A9grale_de_Riemann})
\item Contexte : Intégrale de Lebesgue --- Wikipédia (\url{https://fr.wikipedia.org/wiki/Int%C3%A9grale_de_Lebesgue})
\item Lecture : Rudin, \textit{Principles of Mathematical Analysis}, ch. 11
\item Lecture : Bressoud, \textit{A Radical Approach to Lebesgue's Theory of Integration}
\end{refs}

\bigskip

\begin{problem}[{Construire la mesure de Lebesgue --- Master}]
\textbf{RA-15.} \textbf{Problème.} Définissez la mesure extérieure de $ E \subset \mathbb{R} $ par
$ m^*(E) = \inf \sum |I_k| $, la borne inférieure des longueurs totales sur tous les
recouvrements dénombrables de $ E $ par des intervalles ouverts $ I_k $.
\begin{enumerate}
\item[(a)] Démontrez que $ m^* $ est monotone, dénombrablement
sous-additive, invariante par translation, et qu'elle attribue à chaque intervalle sa longueur.
\item[(b)] Dites que $ E $ est
mesurable (au sens de Carathéodory) si $ m^*(A) = m^*(A \cap E) + m^*(A \setminus E) $ pour tout
ensemble $ A $. Démontrez que les ensembles mesurables forment une tribu contenant tous les
intervalles, et que $ m^* $ restreinte à cette tribu est dénombrablement additive.
\item[(c)] Démontrez que l'ensemble de Cantor est mesurable
et de mesure nulle, et construisez un \og{} ensemble de Cantor gras \fg{} --- fermé, d'intérieur vide, mais de
mesure strictement positive.
\end{enumerate}

\end{problem}

\noindent Borel (1898) et Lebesgue (1902) ont construit la mesure sur la droite ; la
condition de découpage de Carathéodory (1914) a distillé la construction sous la forme aujourd'hui
standard, qui fonctionne mot pour mot dans $ \mathbb{R}^n $ et bien au-delà. L'ensemble de Cantor
gras de la question (c), dû à Smith (1875) et à Volterra, alimente un contre-exemple célèbre : une
fonction dérivable dont la dérivée est bornée mais n'est pas Riemann-intégrable, l'un des irritants
qui ont rendu la théorie de Lebesgue nécessaire.

\begin{refs}
\item Contexte : Mesure de Lebesgue --- Wikipédia (\url{https://fr.wikipedia.org/wiki/Mesure_de_Lebesgue})
\item Contexte : Critère de Carathéodory --- Wikipédia (en anglais) (\url{https://en.wikipedia.org/wiki/Carath%C3%A9odory%27s_criterion})
\item Lecture : Stein \& Shakarchi, \textit{Real Analysis}, ch. 1
\item Cours : MIT OCW 18.125 Measure and Integration (\url{https://ocw.mit.edu/courses/18-125-measure-and-integration-fall-2003/})
\end{refs}

\bigskip

\begin{problem}[{Les théorèmes de convergence à l'œuvre --- Master}]
\textbf{RA-16.} \textbf{Problème.}
\begin{enumerate}
\item[(a)] Démontrez le théorème de convergence monotone : si
$ 0 \le f_1 \le f_2 \le \cdots $ sont mesurables et si $ f_n \to f $ simplement, alors
$ \int f_n \to \int f $.
\item[(b)] Déduisez-en le lemme de Fatou et le théorème de convergence dominée : si
$ f_n \to f $ simplement et $ |f_n| \le g $ pour une fonction intégrable fixée $ g $, alors
$ \int f_n \to \int f $.
\item[(c)] Montrez que l'hypothèse de domination ne peut pas être supprimée, à l'aide de la
bosse qui s'échappe de RA-10(c).
\item[(d)] Application : justifiez soigneusement la limite classique
\[ \int_0^n \left(1 - \frac{x}{n}\right)^n x^{s-1}\,dx \;\longrightarrow\;
\int_0^\infty e^{-x} x^{s-1}\,dx = \Gamma(s) \qquad (s > 0). \]
\end{enumerate}

\end{problem}

\noindent Ces théorèmes --- Beppo Levi (1906), Fatou (1906), Lebesgue (1904--1910) ---
constituent la récompense de toute la construction : l'intégrale de Lebesgue passe à la limite sous
des hypothèses réellement vérifiables, là où l'intégrale de Riemann exige la convergence uniforme.
Le point (d) est le genre d'interversion qu'Euler pratiquait sans crainte en 1729 lorsqu'il a
découvert la fonction Gamma ; la théorie de la mesure est la façon dont nous en payons désormais
honnêtement le prix.

\begin{refs}
\item Contexte : Théorème de convergence dominée --- Wikipédia (\url{https://fr.wikipedia.org/wiki/Th%C3%A9or%C3%A8me_de_convergence_domin%C3%A9e})
\item Contexte : Théorème de convergence monotone --- Wikipédia (\url{https://fr.wikipedia.org/wiki/Th%C3%A9or%C3%A8me_de_convergence_monotone})
\item Lecture : Stein \& Shakarchi, \textit{Real Analysis}, ch. 2
\item Lecture : Folland, \textit{Real Analysis: Modern Techniques and Their Applications}, ch. 2
\end{refs}

\bigskip

\begin{problem}[{Quand Fubini échoue --- Master}]
\textbf{RA-17.} \textbf{Problème.} Soit
$ f(x, y) = \dfrac{x^2 - y^2}{(x^2 + y^2)^2} $ sur le carré $ (0, 1] \times (0, 1] $.
\begin{enumerate}
\item[(a)] Calculez exactement les deux intégrales itérées
$ \int_0^1 \!\! \int_0^1 f \,dx\,dy $ et $ \int_0^1 \!\! \int_0^1 f \,dy\,dx $,
et montrez qu'elles valent $ \pi/4 $ et $ -\pi/4 $ dans un certain ordre : l'ordre
d'intégration compte.
\item[(b)] Énoncez précisément les théorèmes de Fubini et de Tonelli et identifiez l'hypothèse que
$ f $ viole (montrez que $ \iint |f| = \infty $).
\item[(c)] Construisez un second échec sans aucune
astuce de compensation : une fonction sur $ [0,1]^2 $ dont les deux intégrales itérées existent
mais diffèrent, construite à partir de masses placées sur une grille diagonale.
\end{enumerate}

\end{problem}

\noindent Cauchy connaissait déjà de tels exemples en 1827. Fubini (1907) et Tonelli (1909)
ont tracé la bonne frontière : l'intégration itérée est sûre pour les fonctions intégrables, et sans
condition pour les fonctions positives. L'exemple de (a) est la mise en garde classique enseignée à
chaque génération, car dans le travail appliqué l'interversion des intégrales se pratique
quotidiennement et la vérification de l'intégrabilité absolue est la seule chose qui vous sépare de
$ -\pi/4 $ au lieu de $ \pi/4 $.

\begin{refs}
\item Contexte : Théorème de Fubini --- Wikipédia (\url{https://fr.wikipedia.org/wiki/Th%C3%A9or%C3%A8me_de_Fubini})
\item Lecture : Folland, \textit{Real Analysis}, ch. 2
\item Lecture : Stein \& Shakarchi, \textit{Real Analysis}, ch. 2
\end{refs}

\bigskip

\begin{problem}[{Un ensemble qu'on ne peut pas mesurer --- Master}]
\textbf{RA-18.} \textbf{Problème.}
\begin{enumerate}
\item[(a)] Construisez un ensemble de Vitali : en utilisant l'axiome du choix,
choisissez un représentant dans chaque classe d'équivalence de $ \mathbb{R} $ pour la relation
$ x \sim y $ si et seulement si $ x - y \in \mathbb{Q} $, tous les représentants étant pris dans
$ [0, 1] $. Démontrez que l'ensemble $ V $ obtenu n'est pas Lebesgue-mesurable, en montrant
qu'une infinité dénombrable de translatés rationnels de $ V $ pavent un ensemble coincé entre
$ [0, 1] $ et $ [-1, 2] $, et en tirant une contradiction de l'additivité dénombrable et de
l'invariance par translation, que l'on ait $ m(V) = 0 $ ou $ m(V) > 0 $.
\item[(b)] Concluez à la morale plus forte :
il n'existe \textit{aucune} mesure définie sur \textit{toutes} les parties de $ \mathbb{R} $ qui soit
dénombrablement additive, invariante par translation et donne à $ [0, 1] $ la mesure 1.
\end{enumerate}

\end{problem}

\noindent Giuseppe Vitali a publié la construction en 1905, moins de trois ans après la
thèse de Lebesgue : la nouvelle mesure ne pourrait jamais couvrir tous les ensembles. La
démonstration utilise l'axiome du choix de manière essentielle --- Solovay a montré en 1970 que dans
un modèle convenable de la théorie des ensembles sans choix intégral, toute partie de $\mathbb{R}$ est
mesurable. Les ensembles non mesurables ne sont donc pas des objets que l'on puisse exhiber, mais
seulement faire exister par démonstration ; savoir si cela compte comme une existence est une
question avec laquelle les mathématiques ont appris à vivre.

\begin{refs}
\item Contexte : Ensemble de Vitali --- Wikipédia (\url{https://fr.wikipedia.org/wiki/Ensemble_de_Vitali})
\item Contexte : Ensemble non mesurable --- Wikipédia (en anglais) (\url{https://en.wikipedia.org/wiki/Non-measurable_set})
\item Lecture : Royden \& Fitzpatrick, \textit{Real Analysis}
\end{refs}

\bigskip

\begin{problem}[{Banach--Tarski, compris --- Master}]
\textbf{RA-19.} \textbf{Problème.} (Problème de compréhension.) Le paradoxe de Banach--Tarski : une
boule pleine de $ \mathbb{R}^3 $ peut être partagée en un nombre fini de morceaux qui peuvent être
réassemblés, à l'aide de rotations et de translations seulement, en deux boules chacune congruente à
la boule d'origine. Lisez un traitement complet et rédigez un compte rendu répondant à :
\begin{enumerate}
\item[(a)] Pourquoi cela est-il compatible
avec RA-18 --- quels morceaux doivent nécessairement être non mesurables, et pourquoi le
\og{} dédoublement \fg{} ne viole-t-il pas la conservation du volume ?
\item[(b)] Où le groupe libre à deux générateurs se cache-t-il à l'intérieur du
groupe des rotations $ SO(3) $, et comment sa décomposition paradoxale fait-elle marcher la
construction ?
\item[(c)] Pourquoi n'y a-t-il pas de tel paradoxe dans $ \mathbb{R} $ ni dans $ \mathbb{R}^2 $ ---
qu'est-ce qu'un groupe moyennable, et qu'a démontré Banach sur les mesures finiment additives en
petite dimension ?
\end{enumerate}

\end{problem}

\noindent Hausdorff en a trouvé le germe en 1914 ; Banach et Tarski ont publié le paradoxe
complet en 1924. On le raconte souvent comme du mysticisme, mais son contenu réel est un théorème
sur les groupes : les rotations en dimension trois contiennent assez de liberté pour être
paradoxales, tandis que les isométries de la droite et du plan n'en contiennent pas --- observation à
partir de laquelle von Neumann a créé la moyennabilité (1929). C'est la réponse définitive à qui
voudrait que tous les ensembles aient un volume.

\begin{refs}
\item Contexte : Paradoxe de Banach--Tarski --- Wikipédia (\url{https://fr.wikipedia.org/wiki/Paradoxe_de_Banach-Tarski})
\item Lecture : Wagon, \textit{The Banach--Tarski Paradox}
\item Lecture : Tao, \textit{An Introduction to Measure Theory}
\end{refs}

\bigskip

\begin{problem}[{Le théorème de différentiation de Lebesgue --- Master}]
\textbf{RA-20.} \textbf{Problème.} Démontrez le théorème de différentiation de Lebesgue : si $ f $ est
intégrable sur $ \mathbb{R}^n $, alors pour presque tout $ x $ les moyennes de $ f $ sur des
boules $ B $ se contractant vers $ x $ convergent vers $ f(x) $ --- en symboles,
$ \frac{1}{|B|} \int_B f \to f(x) $. Suivez la route standard et démontrez chaque ingrédient :
\begin{enumerate}
\item[(a)] le lemme de recouvrement de Vitali ;
\item[(b)] l'inégalité maximale de Hardy--Littlewood --- la fonction
maximale $ Mf(x) = \sup_{B \ni x} \frac{1}{|B|} \int_B |f| $ vérifie la majoration de type faible
$ |\{ Mf > \lambda \}| \le C \, \|f\|_1 / \lambda $ ;
\item[(c)] le passage de l'inégalité maximale et de la
densité des fonctions continues à l'énoncé \og{} presque partout \fg{}. Déduisez-en le théorème fondamental
de l'analyse à une variable pour les intégrales de Lebesgue : l'intégrale indéfinie d'une fonction
intégrable est dérivable presque partout, de dérivée $ f $.
\end{enumerate}

\end{problem}

\noindent Lebesgue avait démontré le théorème en dimension un dès 1910 ; Hardy et
Littlewood ont introduit la fonction maximale en 1930, dans un article au parfum de cricket sur les
moyennes de batteurs, et leur inégalité est devenue le modèle de centaines de résultats \og{} presque
partout \fg{}. Ce théorème est le théorème fondamental de l'analyse version théorie de la mesure, et son
schéma de démonstration --- lemme de recouvrement, inégalité maximale, sous-classe dense --- est la
porte d'entrée de l'analyse harmonique et des problèmes de recherche ci-dessous.

\begin{refs}
\item Contexte : Théorème de différentiation de Lebesgue --- Wikipédia (\url{https://fr.wikipedia.org/wiki/Th%C3%A9or%C3%A8me_de_diff%C3%A9rentiation_de_Lebesgue})
\item Contexte : Fonction maximale de Hardy--Littlewood --- Wikipédia (\url{https://fr.wikipedia.org/wiki/Fonction_maximale_de_Hardy-Littlewood})
\item Lecture : Stein \& Shakarchi, \textit{Real Analysis}, ch. 3
\item Cours : MIT OCW 18.125 Measure and Integration (\url{https://ocw.mit.edu/courses/18-125-measure-and-integration-fall-2003/})
\end{refs}

\bigskip

\begin{problem}[{Stone--Weierstrass --- Master, short}]
\textbf{RA-30.} \textbf{Problème.} Soit $ K $ un espace compact séparé et soit
$ \mathcal{A} \subseteq C(K, \mathbb{R}) $ une sous-algèbre contenant les fonctions constantes et
séparant les points (pour $ x \ne y $ il existe $ g \in \mathcal{A} $ avec
$ g(x) \ne g(y) $). Démontrez que $ \mathcal{A} $ est dense dans $ C(K, \mathbb{R}) $ pour la
norme uniforme, et déduisez-en le théorème de Weierstrass selon lequel toute fonction continue sur
$ [0,1] $ est limite uniforme de polynômes.
\end{problem}

\noindent Weierstrass a démontré le cas polynomial en 1885, à soixante-dix ans, comme une
sorte de contrepoids au monstre de RA-11 : les fonctions continues peuvent être pathologiques, et
pourtant chacune d'elles est approchable par les fonctions les plus dociles qui soient. La
généralisation de Marshall Stone, en 1937, a dépouillé la démonstration jusqu'à ses os algébriques ---
une algèbre de fonctions, stable par produit, capable de distinguer les points est déjà tout ce
qu'il faut. La démonstration de Bernstein du cas classique (1912), qui utilise la loi des grands
nombres, mérite d'être lue en parallèle ; et la version complexe est fausse sans stabilité par
conjugaison, puisque les limites uniformes de polynômes en $ z $ sont holomorphes.

\begin{refs}
\item Contexte : Théorème de Stone--Weierstrass --- Wikipédia (\url{https://fr.wikipedia.org/wiki/Th%C3%A9or%C3%A8me_de_Stone-Weierstrass})
\item Lecture : Rudin, \textit{Principles of Mathematical Analysis}, ch. 7
\item Lecture : Folland, \textit{Real Analysis}, 2e éd., ch. 4
\end{refs}

\bigskip

\begin{problem}[{Le théorème d'Egorov --- Master, short}]
\textbf{RA-31.} \textbf{Problème.} Soit $ (X, \mathcal{M}, \mu) $ un espace mesuré avec
$ \mu(X) < \infty $, et soient des fonctions mesurables $ f_n $ convergeant simplement presque
partout vers $ f $, toutes finies presque partout. Démontrez le théorème d'Egorov : pour tout
$ \varepsilon > 0 $ il existe un ensemble mesurable $ E \subseteq X $ avec
$ \mu(X \setminus E) < \varepsilon $ sur lequel $ f_n \to f $ uniformément ; montrez ensuite
par un exemple sur $ \mathbb{R} $ que l'hypothèse $ \mu(X) < \infty $ ne peut pas être
supprimée.
\end{problem}

\noindent C'est le deuxième des trois principes de Littlewood --- toute suite convergente de
fonctions mesurables est presque uniformément convergente ---, et il rend parfaitement précise
l'intuition vague selon laquelle la théorie de la mesure est \og{} de l'analyse à un petit ensemble
près \fg{}. Dmitri Egorov l'a publié en 1911 ; Carlo Severini l'avait démontré indépendamment en 1910,
dans une revue italienne, et fut ignoré pendant des décennies. Le théorème est le pont standard des
hypothèses \og{} presque partout \fg{} vers les estimations uniformes, et son échec sur les espaces de
mesure infinie est la raison pour laquelle tant d'énoncés portent une hypothèse de finitude.

\begin{refs}
\item Contexte : Théorème d'Egoroff --- Wikipédia (\url{https://fr.wikipedia.org/wiki/Th%C3%A9or%C3%A8me_d'Egoroff})
\item Lecture : Stein \& Shakarchi, \textit{Real Analysis}, ch. 1
\item Lecture : Folland, \textit{Real Analysis}, 2e éd., ch. 2
\end{refs}

\bigskip

\begin{problem}[{Arzelà--Ascoli : la compacité pour les fonctions --- Master}]
\textbf{RA-32.} \textbf{Problème.} Soit $ K $ un espace métrique compact. Une famille $ \mathcal{F} $ de
fonctions $ K \to \mathbb{R} $ est \textit{équicontinue} si pour tout
$ \varepsilon > 0 $ il existe un $ \delta > 0 $, le même pour tout
$ f \in \mathcal{F} $, tel que $ d(x,y) < \delta $ entraîne
$ |f(x) - f(y)| < \varepsilon $.
\begin{enumerate}
\item[(a)] Démontrez le théorème d'Arzelà--Ascoli : une partie de $ C(K) $ est d'adhérence compacte pour
la norme uniforme si et seulement si elle est bornée et équicontinue. (Les deux sens ; le sens
substantiel est l'extraction diagonale d'une sous-suite.)
\item[(b)] Montrez qu'aucune des deux hypothèses ne peut être supprimée : exhibez une suite bornée de
$ C[0,1] $ sans sous-suite uniformément convergente, et une suite équicontinue qui n'en a pas non
plus.
\item[(c)] Déduisez-en qu'une suite de fonctions partageant une même constante de Lipschitz et une même
borne uniforme admet une sous-suite uniformément convergente, et servez-vous-en pour démontrer le
théorème d'existence de Peano : si $ F $ est continue au voisinage de $ (t_0, y_0) $, alors
$ y' = F(t, y) $, $ y(t_0) = y_0 $ admet une solution sur un intervalle autour de $ t_0 $ ---
sans aucune condition de Lipschitz. (Construisez des solutions approchées par lignes polygonales
d'Euler et extrayez une limite.) Montrez par un exemple que l'unicité est réellement mise en
défaut.
\item[(d)] Expliquez où le théorème tombe en défaut pour $ C(\mathbb{R}) $ muni de la norme uniforme, et
comment l'énoncé se répare en passant à la convergence uniforme locale.
\end{enumerate}

\end{problem}

\noindent Giulio Ascoli (1883--1884) et Cesare Arzelà (1889, 1895) cherchaient à rendre
rigoureux le genre d'argument que Riemann avait employé sans précaution dans le principe de
Dirichlet : prendre une famille minimisante de courbes et passer à la limite. Leur théorème est le
premier théorème de compacité en dimension infinie, et tous ceux qui ont suivi --- Rellich--Kondrachov
pour les espaces de Sobolev, Montel pour les familles holomorphes, Prokhorov pour les mesures de
probabilité --- reposent sur la même intuition : un module de continuité uniforme se substitue à la
dimension finie. Le théorème d'existence de Peano (1890) en est la retombée la plus célèbre :
abandonner l'hypothèse de Lipschitz achète l'existence au prix de l'unicité, ce qui est l'état
honnête des choses pour un très grand nombre d'équations différentielles.

\begin{refs}
\item Contexte : Théorème d'Ascoli --- Wikipédia (\url{https://fr.wikipedia.org/wiki/Th%C3%A9or%C3%A8me_d'Ascoli})
\item Contexte : Équicontinuité --- Wikipédia (\url{https://fr.wikipedia.org/wiki/%C3%89quicontinuit%C3%A9})
\item Contexte : Théorème de Cauchy--Peano--Arzelà --- Wikipédia (\url{https://fr.wikipedia.org/wiki/Th%C3%A9or%C3%A8me_de_Cauchy-Peano-Arzel%C3%A0})
\item Lecture : Rudin, \textit{Principles of Mathematical Analysis}, ch. 7
\end{refs}

\bigskip

\begin{problem}[{Continuité absolue, et le vrai théorème fondamental --- Master}]
\textbf{RA-40.} \textbf{Problème.} Dites que $ f : [a,b] \to \mathbb{R} $ est \textit{absolument continue} si pour
tout $ \varepsilon > 0 $ il existe un $ \delta > 0 $ tel que pour toute famille finie
d'intervalles deux à deux disjoints $ (a_k,b_k) \subseteq [a,b] $ avec
$ \sum_k (b_k - a_k) < \delta $ on ait $ \sum_k |f(b_k)-f(a_k)| < \varepsilon $.
\begin{enumerate}
\item[(a)] Démontrez la chaîne d'implications : lipschitzienne $ \Rightarrow $ absolument continue
$ \Rightarrow $ uniformément continue et à variation bornée (RA-29). Montrez que chaque
implication est stricte --- $ \sqrt{x} $ sur $ [0,1] $ est absolument continue sans être
lipschitzienne, la fonction de Cantor de RA-38 est continue et à variation bornée sans être
absolument continue, et $ x\sin(1/x) $ est uniformément continue sans être à variation
bornée.
\item[(b)] Démontrez le théorème fondamental de l'analyse pour l'intégrale de Lebesgue : $ f $ est
absolument continue sur $ [a,b] $ si et seulement si $ f $ est dérivable presque partout,
$ f' $ est intégrable, et
\[ f(x) - f(a) = \int_a^x f'(t)\,dt \quad \text{pour tout } x \in [a,b]. \]
(Un sens est le théorème de différentiation de Lebesgue de RA-20 ; l'autre demande un argument de
recouvrement de Vitali.)
\item[(c)] Démontrez que $ f $ est lipschitzienne de constante $ M $ si et seulement si elle est
absolument continue avec $ |f'| \le M $ presque partout.
\item[(d)] Démontrez que les fonctions absolument continues ont la propriété (N) de Lusin : elles envoient
les ensembles négligeables sur des ensembles négligeables. Montrez que la fonction de Cantor ne l'a
pas, puis démontrez le théorème de Banach--Zarecki : $ f $ est absolument continue si et
seulement si elle est continue, à variation bornée, et possède la propriété (N).
\end{enumerate}

\end{problem}

\noindent Giuseppe Vitali a isolé la continuité absolue en 1905 dans un seul but : décrire
exactement quelles fonctions sont des intégrales de Lebesgue indéfinies. C'est la réponse définitive
à la question qu'ouvre CAL-09(iii) et qu'affine RA-38 --- \textit{quelles} fonctions sont récupérées en
intégrant leur dérivée --- et la réponse n'est ni \og{} les fonctions dérivables \fg{}, ni \og{} les fonctions
continues monotones \fg{}, mais précisément cette classe-là. Lebesgue avait démontré en 1904 qu'une
fonction monotone est dérivable presque partout ; la question (b) est ce à quoi ce théorème sert, et
la question (d), due à Banach et Zarecki vers 1925--1930, décompose la condition en trois
morceaux vérifiables indépendamment.

\begin{refs}
\item Contexte : Absolue continuité --- Wikipédia (\url{https://fr.wikipedia.org/wiki/Absolue_continuit%C3%A9})
\item Contexte : Application lipschitzienne --- Wikipédia (\url{https://fr.wikipedia.org/wiki/Application_lipschitzienne})
\item Lecture : Royden \& Fitzpatrick, \textit{Real Analysis}, 4e éd., ch. 6
\item Lecture : Wheeden \& Zygmund, \textit{Measure and Integral} (le chapitre sur la différentiation)
\end{refs}

\bigskip

\begin{problem}[{Le théorème de Lusin --- Master, short}]
\textbf{RA-41.} \textbf{Problème.} Démontrez le théorème de Lusin : si $ f $ est mesurable et
finie presque partout sur un ensemble $ E \subseteq \mathbb{R} $ de mesure finie, alors pour tout
$ \varepsilon > 0 $ il existe un fermé $ K \subseteq E $ avec
$ m(E \setminus K) < \varepsilon $ tel que la restriction $ f|_K $ soit continue. Montrez
ensuite le peu que cela dit sur $ f $ elle-même : appliquez-le explicitement à la fonction de
Dirichlet de RA-05, qui n'est continue en absolument aucun point de $ \mathbb{R} $.
\end{problem}

\noindent C'est le troisième des trois principes de Littlewood --- toute fonction mesurable
est presque continue --- et la distinction que la seconde tâche impose, entre \og{} $ f $ restreinte à
$ K $ est continue \fg{} et \og{} $ f $ est continue aux points de $ K $ \fg{}, est exactement le piège
que le théorème tend. Nikolaï Luzin l'a démontré en 1912, une version se trouvant déjà chez Vitali
(1905) ; joint au théorème d'Egorov (RA-31) et à l'approximation des ensembles mesurables par des
réunions finies d'intervalles (RA-15), il rend précise l'idée que la théorie de la mesure est de
l'analyse ordinaire pratiquée à un ensemble exceptionnel arbitrairement petit près.

\begin{refs}
\item Contexte : Théorème de Lusin --- Wikipédia (\url{https://fr.wikipedia.org/wiki/Th%C3%A9or%C3%A8me_de_Lusin})
\item Contexte : Les trois principes de Littlewood en analyse réelle --- Wikipédia (en anglais) (\url{https://en.wikipedia.org/wiki/Littlewood%27s_three_principles_of_real_analysis})
\item Lecture : Stein \& Shakarchi, \textit{Real Analysis}, ch. 1
\end{refs}

\bigskip

\begin{problem}[{Ce qu'est vraiment une forme linéaire sur C[0,1] --- Master}]
\textbf{RA-42.} \textbf{Problème.} Munissez $ C[0,1] $ de la norme de la borne supérieure et soit $ L $ une
forme linéaire continue sur cet espace, c'est-à-dire une application linéaire
$ L : C[0,1] \to \mathbb{R} $ telle que $ |L(f)| \le M \|f\|_{\infty} $ pour une certaine
constante $ M $.
\begin{enumerate}
\item[(a)] Démontrez le théorème de F. Riesz de 1909 : il existe une fonction $ g $ à variation bornée
sur $ [0,1] $ telle que
\[ L(f) = \int_0^1 f \, dg \qquad \text{pour tout } f \in C[0,1], \]
l'intégrale étant une intégrale de Riemann--Stieltjes, et $ \|L\| = V_0^1(g) $ lorsque
$ g $ est normalisée (disons $ g(0)=0 $ et $ g $ continue à droite sur $ (0,1) $). Montrez
par un exemple que sans normalisation $ g $ n'est pas unique.
\item[(b)] Démontrez le cas positif sous forme mesurée : si $ L(f) \ge 0 $ dès que $ f \ge 0 $, alors
il existe une mesure borélienne finie $ \mu $ sur $ [0,1] $ telle que $ L(f) = \int f \, d\mu $
pour toute $ f $ continue, et $ \mu $ est unique.
\item[(c)] Déduisez-en le théorème de sélection de Helly, et avec lui la compacité faible-* de l'ensemble
des mesures boréliennes de probabilité sur $ [0,1] $ : toute suite de telles mesures admet une
sous-suite $ \mu_{n_k} $ et une limite $ \mu $ avec $ \int f d\mu_{n_k} \to \int f d\mu $ pour
toute $ f $ continue. Expliquez le rapport avec le théorème de Banach--Alaoglu, et montrez
que l'énoncé correspondant pour la topologie de la norme est faux.
\item[(d)] Montrez les dents du théorème : démontrez qu'aucune fonction intégrable $ h \in L^1[0,1] $ ne
vérifie $ f(0) = \int_0^1 f h $ pour toute $ f $ continue. La forme linéaire d'évaluation est
donc une mesure et non une fonction --- c'est le point où la distribution de Dirac cesse d'être un
abus de notation et devient un objet.
\end{enumerate}

\end{problem}

\noindent Frigyes Riesz a annoncé ce résultat dans une note de quatre pages aux
\textit{Comptes rendus} en 1909, et c'est sans doute là que commence l'analyse fonctionnelle : pour la
première fois un dual était identifié concrètement, et la réponse n'était pas un espace de fonctions
mais un espace de mesures. Markov et Kakutani l'ont étendu aux espaces compacts séparés vers
1938--1941, et l'énoncé moderne porte les trois noms. La question (d) est le germe historique de
la théorie des distributions : la construction de Schwartz dans les années 1940 est la réponse
systématique à la découverte que les objets contre lesquels les analystes veulent intégrer sont plus
généraux que les fonctions dont ils étaient partis.

\begin{refs}
\item Contexte : Théorème de représentation de Riesz (Fréchet--Riesz) --- Wikipédia (\url{https://fr.wikipedia.org/wiki/Th%C3%A9or%C3%A8me_de_repr%C3%A9sentation_de_Riesz_(Fr%C3%A9chet-Riesz)})
\item Contexte : Théorème de représentation de Riesz--Markov--Kakutani --- Wikipédia (\url{https://fr.wikipedia.org/wiki/Th%C3%A9or%C3%A8me_de_repr%C3%A9sentation_de_Riesz_(Riesz-Markov)})
\item Contexte : Intégrale de Riemann--Stieltjes --- Wikipédia (\url{https://fr.wikipedia.org/wiki/Int%C3%A9grale_de_Stieltjes})
\item Lecture : Rudin, \textit{Real and Complex Analysis}, ch. 2
\item Lecture : Folland, \textit{Real Analysis}, 2e éd., ch. 7 (mesures de Radon)
\end{refs}

\bigskip

\begin{problem}[{Steinhaus : une mesure strictement positive fait de la place --- Master}]
\textbf{RA-43.} \textbf{Problème.} Soit $ E \subseteq \mathbb{R} $ Lebesgue-mesurable avec
$ m(E) > 0 $, et soit $ E - E = \{ x - y : x, y \in E \} $.
\begin{enumerate}
\item[(a)] Démontrez le théorème de Steinhaus : $ E - E $ contient un intervalle ouvert autour de
$ 0 $. Donnez deux démonstrations --- l'une à partir du théorème de densité de Lebesgue, en prenant
un point où $ E $ occupe au moins 99 \% d'un petit intervalle ; l'autre en démontrant que
$ \varphi(t) = m\big( E \cap (E+t) \big) $ est continue en $ t $ et strictement positive en
$ t = 0 $.
\item[(b)] Déduisez-en qu'un sous-groupe mesurable de $ (\mathbb{R},+) $ de mesure strictement positive
est $ \mathbb{R} $ tout entier, et qu'un sous-groupe mesurable de mesure nulle est un sous-groupe
strict d'intérieur vide --- les sous-groupes mesurables sont donc exactement de deux sortes.
\item[(c)] Démontrez que toute solution Lebesgue-mesurable de l'équation fonctionnelle de Cauchy
$ f(x+y) = f(x)+f(y) $ est linéaire, $ f(x) = cx $. Expliquez ensuite, en vous référant à
RA-18, pourquoi les solutions non linéaires construites à partir d'une base de Hamel de
$ \mathbb{R} $ sur $ \mathbb{Q} $ ne peuvent jamais être exhibées : leurs graphes sont denses
dans le plan, et aucune d'elles n'est mesurable.
\item[(d)] Montrez que ni l'hypothèse ni la conclusion ne peuvent être renversées : l'ensemble de Cantor
gras de RA-15(c) est de mesure strictement positive et ne contient aucun intervalle, donc (a) porte
bien sur l'ensemble des différences ; et l'ensemble triadique de Cantor $ C $ est de mesure nulle
alors que $ C - C = [-1,1] $, donc la réciproque de (a) est fausse.
\end{enumerate}

\end{problem}

\noindent Hugo Steinhaus a publié cela en 1920, dans le tout premier volume de
\textit{Fundamenta Mathematicae}, la revue que l'école polonaise venait de fonder. Le théorème est le
pont standard de la mesure vers l'algèbre, et la raison pour laquelle tant d'énoncés de rigidité
commencent par \og{} supposons l'application mesurable \fg{} : la mesurabilité est une hypothèse d'apparence
anodine qui interdit secrètement toute pathologie, comme le montre (c) pour les fonctions additives
que Cauchy avait étudiées en 1821. André Weil a démontré l'analogue dans les groupes localement
compacts --- pour un ensemble de mesure de Haar strictement positive, $ EE^{-1} $ est un voisinage
de l'élément neutre ---, et c'est le moteur derrière la continuité automatique des homomorphismes
mesurables.

\begin{refs}
\item Contexte : Théorème de Steinhaus --- Wikipédia (\url{https://fr.wikipedia.org/wiki/Th%C3%A9or%C3%A8me_de_Steinhaus})
\item Contexte : Équation fonctionnelle de Cauchy --- Wikipédia (\url{https://fr.wikipedia.org/wiki/%C3%89quation_fonctionnelle_de_Cauchy})
\item Contexte : Théorème de densité de Lebesgue --- Wikipédia (en anglais) (\url{https://en.wikipedia.org/wiki/Lebesgue%27s_density_theorem})
\item Article : K. Stromberg, \og{} An elementary proof of Steinhaus's theorem \fg{}, \textit{Proc. Amer. Math. Soc.} 36 (1972), 308
\item Lecture : Oxtoby, \textit{Measure and Category}
\end{refs}

\bigskip


\section*{Recherche}

\begin{problem}[{Le problème de Kakeya --- Recherche}]
\textbf{RA-21.} \textbf{Problème.} Un \textit{ensemble de Kakeya} (ensemble de Besicovitch) dans
$ \mathbb{R}^n $ est un compact contenant un segment de longueur 1 dans chaque direction.
\begin{enumerate}
\item[(a)] Échauffement, résolu depuis longtemps : lisez la construction de Besicovitch montrant qu'un
ensemble de Kakeya du plan peut être de mesure de Lebesgue nulle.
\item[(b)] Démontrez la conjecture de Kakeya en dimension deux
(Davies, 1971) : tout ensemble de Kakeya de $ \mathbb{R}^2 $ est de dimension de Hausdorff 2.
\item[(c)] La
conjecture des ensembles de Kakeya affirme que tout ensemble de Kakeya de $ \mathbb{R}^n $ est de
dimension de Hausdorff et de dimension de Minkowski $ n $. Étudiez la démonstration de Wang--Zahl
pour $ n = 3 $ à travers l'introduction expositive de Guth, et rendez compte des notions clés de
configurations de tubes granuleuses (\textit{grainy}) et collantes (\textit{sticky}). Le cas
$ n \ge 4 $ est ouvert : démontrez la conjecture en une dimension $ \ge 4 $ quelconque.
\end{enumerate}

\end{problem}

\noindent Kakeya a demandé en 1917 quelle est la plus petite aire dans laquelle on peut
retourner une aiguille de longueur 1 ; la réponse de Besicovitch (1919, 1928) --- une aire
arbitrairement petite --- a créé la question de la dimension, devenue un pilier de l'analyse
harmonique, relié à la restriction (RA-22), aux EDP et à la combinatoire. \textbf{Statut :} les
dimensions 1 et 2 sont classiques ; en février 2025, Hong Wang et Joshua Zahl ont annoncé une
démonstration de la conjecture des ensembles de Kakeya en dimension trois dans une prépublication
longue et complexe. En 2026, la démonstration a été étudiée de près et est acceptée par les
spécialistes du domaine --- il existe des exposés de Guth et de Tao --- et Wang a reçu la médaille
Fields 2026 (annoncée au CIM de Philadelphie en juillet 2026) en grande partie pour ce travail,
devenant la troisième femme à l'obtenir. La conjecture reste ouverte en toute dimension
$ n \ge 4 $.

\begin{refs}
\item Contexte : Ensemble de Besicovitch (ensemble de Kakeya) --- Wikipédia (\url{https://fr.wikipedia.org/wiki/Ensemble_de_Besicovitch})
\item Article : Wang \& Zahl, \og{} Volume estimates for unions of convex sets, and the Kakeya set conjecture in three dimensions \fg{} (2025), arXiv:2502.17655 (\url{https://arxiv.org/abs/2502.17655})
\item Article : Guth, \og{} Introduction to the proof of the Kakeya conjecture \fg{} (2025), arXiv:2505.07695 (\url{https://arxiv.org/abs/2505.07695})
\item Lecture : Tao, \og{} The three-dimensional Kakeya conjecture, after Wang and Zahl \fg{} (blog, 2025) (\url{https://terrytao.wordpress.com/2025/02/25/the-three-dimensional-kakeya-conjecture-after-wang-and-zahl/})
\end{refs}

\bigskip

\begin{problem}[{La conjecture de restriction --- Recherche}]
\textbf{RA-22.} \textbf{Problème.} La transformée de Fourier d'une fonction $ L^1 $ est
continue, on peut donc la restreindre à une sphère ; pour $ L^2 $ elle n'est définie que presque
partout, et une sphère est de mesure nulle. La conjecture de restriction de Stein situe la
frontière : pour la sphère unité $ S^{n-1} \subset \mathbb{R}^n $, la restriction à la sphère de
la transformée de Fourier de $ f $ devrait vérifier
\[ \| \hat{f} \|_{L^q(S^{n-1})} \le C \, \| f \|_{L^p(\mathbb{R}^n)}
\qquad \text{pour tout } p < \tfrac{2n}{n+1} \]
(avec $ q $ déterminé par l'homogénéité) --- c'est la courbure de la sphère qui rend possible une
telle estimation.
\begin{enumerate}
\item[(a)] Démontrez le théorème de Tomas--Stein, le résultat partiel fondamental,
dans le domaine qu'il couvre.
\item[(b)] Lisez comment une estimation de restriction complète impliquerait la conjecture de Kakeya de
RA-21.
\item[(c)] Démontrez la conjecture en une dimension $ n \ge 3 $ quelconque.
\end{enumerate}

\end{problem}

\noindent Stein a observé à la fin des années 1960 que la courbure permet de restreindre
la transformée de Fourier à des surfaces, lançant l'analyse harmonique euclidienne moderne ; le
théorème du multiplicateur de la boule de Fefferman (1971) a révélé le lien profond avec les
ensembles de Kakeya. \textbf{Statut :} le cas de la dimension deux a été réglé dans les années 1970
(Fefferman, Zygmund) ; pour tout $ n \ge 3 $ la conjecture est ouverte en 2026 --- le théorème de
Wang--Zahl de RA-21 démontre une conséquence géométrique nécessaire dans $ \mathbb{R}^3 $, et non
l'estimation de restriction elle-même, et Wang a donné des exposés sur les stratégies d'attaque de
la restriction au-delà de Kakeya. Les méthodes de découplage de Bourgain, Demeter et Guth ont porté
les meilleurs exposants récents.

\begin{refs}
\item Contexte : Conjecture de restriction --- Wikipédia (en anglais) (\url{https://en.wikipedia.org/wiki/Restriction_conjecture})
\item Article : Tao, \og{} Recent progress on the restriction conjecture \fg{} (2003), arXiv:math/0311181 (\url{https://arxiv.org/abs/math/0311181})
\item Lecture : Stein \& Shakarchi, \textit{Functional Analysis} (Princeton Lectures in Analysis IV)
\end{refs}

\bigskip

\begin{problem}[{Le problème des distances de Falconer --- Recherche}]
\textbf{RA-23.} \textbf{Problème.} Pour un ensemble $ E \subset \mathbb{R}^n $, son ensemble
des distances est $ \Delta(E) = \{ |x - y| : x, y \in E \} $. Conjecture de Falconer : si $ E $
est compact de dimension de Hausdorff strictement supérieure à $ n/2 $, alors $ \Delta(E) $ est
de mesure de Lebesgue strictement positive. Démontrez-la --- dans le plan, premier cas ouvert, il faut
montrer qu'une dimension $ > 1 $ force un ensemble des distances de mesure strictement
positive. En guise de porte d'entrée, démontrez la cousine discrète accessible avec des outils
élémentaires : pour $ N $ points du plan, obtenez une minoration non triviale du nombre de
distances deux à deux distinctes (problème d'Erdős de 1946, essentiellement résolu par Guth et Katz
en 2015).
\end{problem}

\noindent Kenneth Falconer a posé la conjecture en 1985 comme analogue continu du problème
des distances distinctes d'Erdős. \textbf{Statut : ouvert en toute dimension.} Le meilleur seuil connu
dans le plan est $ 5/4 $ : Guth, Iosevich, Ou et Wang ont démontré en 2020 que les ensembles plans
de dimension strictement supérieure à $ 5/4 $ ont un ensemble des distances de mesure strictement
positive, en utilisant le même cercle d'idées --- découplage, récurrence sur les échelles --- que RA-21
et RA-22, avec depuis d'autres avancées en dimension supérieure. C'est un problème test permanent
pour mesurer notre compréhension de l'interaction entre mesure, dimension et courbure.

\begin{refs}
\item Contexte : Conjecture de Falconer --- Wikipédia (en anglais) (\url{https://en.wikipedia.org/wiki/Falconer%27s_conjecture})
\item Article : Guth, Iosevich, Ou \& Wang, \og{} On Falconer's distance set problem in the plane \fg{}, \textit{Inventiones Mathematicae} 219 (2020)
\item Lecture : Falconer, \textit{The Geometry of Fractal Sets}
\end{refs}

\bigskip

\begin{problem}[{La conjecture des ensembles spectraux de Fuglede --- Recherche, short}]
\textbf{RA-33.} \textbf{Problème.} Dites qu'un ensemble mesurable borné
$ \Omega \subset \mathbb{R}^d $ de mesure strictement positive est \textit{spectral} si
$ L^2(\Omega) $ admet une base orthogonale d'exponentielles
$ e^{2\pi i \langle \lambda, x \rangle} $ ; Fuglede a conjecturé en 1974 que $ \Omega $ est
spectral si et seulement si des translatés de $ \Omega $ pavent $ \mathbb{R}^d $. Réglez la
conjecture pour $ d = 1 $ et $ d = 2 $ --- \textbf{statut : ouvert dans ces dimensions}, alors
qu'elle est fausse en toute dimension $ d \ge 3 $ et vraie pour tous les corps convexes en toute
dimension.
\end{problem}

\noindent Bent Fuglede est arrivé à la question par la théorie des opérateurs : il se
demandait quand les dérivées partielles sur un domaine admettent des extensions autoadjointes qui
commutent, et a trouvé que la réponse est exactement cette condition géométrique. La conjecture a
tenu trente ans, jusqu'à ce que Terence Tao en réfute un sens en dimension 5 en 2004, à l'aide d'une
construction sur un corps fini transplantée dans $ \mathbb{R}^d $ ; Kolountzakis et Matolcsi ont
ensuite fait descendre les contre-exemples jusqu'à la dimension 3 et tué la réciproque également. La
dimension 1 équivaut à un énoncé purement combinatoire sur les pavages de $ \mathbb{Z} $ --- les
conditions de Coven--Meyerowitz --- et c'est là qu'elle se trouve aujourd'hui, obstinément. Le résultat
positif pour les corps convexes, dû à Lev et Matolcsi en 2022, montre que la conjecture avait raison
exactement pour les ensembles que Fuglede avait en tête.

\begin{refs}
\item Contexte : Conjecture de Fuglede --- Wikipédia (en anglais) (\url{https://en.wikipedia.org/wiki/Fuglede%27s_conjecture})
\item Article : T. Tao, \og{} Fuglede's conjecture is false in 5 and higher dimensions \fg{} (2004), Math. Res. Lett. 11, arXiv:math/0306134 (\url{https://arxiv.org/abs/math/0306134})
\item Article : N. Lev \& M. Matolcsi, \og{} The Fuglede conjecture for convex domains is true in all dimensions \fg{} (2022), Acta Mathematica 228, 385--420, arXiv:1904.12262 (\url{https://arxiv.org/abs/1904.12262})
\end{refs}

\bigskip

\begin{problem}[{La conjecture de Bochner--Riesz --- Recherche, short}]
\textbf{RA-34.} \textbf{Problème.} Pour $ \delta > 0 $, l'opérateur de Bochner--Riesz sur
$ \mathbb{R}^n $ est défini du côté de Fourier par
\[ \widehat{S^{\delta} f}(\xi) \;=\; \left( 1 - |\xi|^2 \right)_{+}^{\delta} \, \hat{f}(\xi), \]
et la conjecture de Bochner--Riesz affirme que $ S^{\delta} $ est borné sur
$ L^p(\mathbb{R}^n) $ exactement lorsque
$ \delta > \max \left\{ n \left| \tfrac{1}{p} - \tfrac{1}{2} \right| - \tfrac{1}{2},\, 0
\right\} $. Démontrez-la en une dimension $ n \ge 3 $ --- \textbf{statut : ouvert pour tout
$ n \ge 3 $}, et réglé dans le plan par Carleson et Sjölin en 1972.
\end{problem}

\noindent Bochner a introduit ces troncatures adoucies dans les années 1930 pour sommer des
intégrales de Fourier qui divergent si l'on coupe simplement au bord d'une sphère. Que la coupure
naïve $ \delta = 0 $ soit réellement sans espoir fut le spectaculaire théorème de Charles Fefferman
de 1971 : le multiplicateur de la boule n'est borné sur $ L^p $ pour aucun $ p \ne 2 $,
démontré en construisant une configuration de tubes fins de type Kakeya. Cet argument a soudé ce
problème à RA-21 et RA-22 --- un théorème de Bochner--Riesz complet impliquerait la conjecture de
restriction, qui impliquerait la conjecture de Kakeya --- et les meilleurs domaines connus en
dimension supérieure proviennent de la machinerie de découplage et de la méthode polynomiale de
Bourgain--Guth et de Guth--Hickman--Iliopoulou.

\begin{refs}
\item Contexte : Moyenne de Bochner--Riesz --- Wikipédia (en anglais) (\url{https://en.wikipedia.org/wiki/Bochner%E2%80%93Riesz_mean})
\item Article : C. Fefferman, \og{} The multiplier problem for the ball \fg{} (1971), Annals of Mathematics 94, 330--336
\item Article : L. Carleson \& P. Sjölin, \og{} Oscillatory integrals and a multiplier problem for the disc \fg{} (1972), Studia Mathematica 44, 287--299
\end{refs}

\bigskip

\begin{problem}[{Le problème de similarité d'Erdős --- Recherche}]
\textbf{RA-35.} \textbf{Problème.} Dites qu'un ensemble $ A \subseteq \mathbb{R} $ est \textit{universel} si tout
ensemble mesurable $ E \subseteq \mathbb{R} $ de mesure de Lebesgue strictement positive contient
une copie affine $ x + tA $ de $ A $, pour un certain $ x \in \mathbb{R} $ et un certain
$ t \ne 0 $. Erdős a demandé vers 1974 si un ensemble infini quelconque peut être universel, et a
conjecturé qu'aucun ne peut l'être.
\begin{enumerate}
\item[(a)] Démontrez que tout ensemble fini est universel. (Le théorème de densité de Lebesgue fait tout le
travail.)
\item[(b)] Démontrez qu'un ensemble non borné n'est jamais universel, de sorte que la question porte sur
les ensembles infinis bornés ; déduisez-en qu'après changement d'échelle il suffit de trancher le
cas d'une suite strictement décroissante $ a_n \to 0 $.
\item[(c)] Lisez et rédigez le théorème de Falconer de 1984 : si $ a_n \downarrow 0 $ lentement, au sens
où $ a_{n+1} / a_n \to 1 $, alors $ \{ a_n \} $ n'est pas universel. Repérez précisément où
l'hypothèse de décroissance lente est utilisée.
\item[(d)] Décidez si la suite géométrique $ \{ 2^{-n} : n \ge 1 \} $ est universelle.
\textbf{Statut : ouvert.} C'est le cas à décroissance rapide le plus simple, et c'est exactement là que
cinquante ans d'efforts se sont enlisés ; la conjecture reste ouverte pour toutes les suites
rapidement décroissantes et pour les ensembles de type Cantor de dimension de Hausdorff nulle.
\end{enumerate}

\end{problem}

\noindent Le problème est un spécimen parfait du goût d'Erdős : un étudiant de première
année comprend l'énoncé, la question (a) tient en trois lignes, et la question (d) a mis tout le
monde en échec pendant un demi-siècle. L'intuition --- un ensemble clairsemé devrait pouvoir être
évité par un ensemble gras --- n'a jamais pu être mise au travail dès que la suite décroît vite, parce
que les constructions naturelles d'ensembles évitants perdent toute leur marge. Les progrès sont
venus surtout en changeant la question : Falconer et Eigen ont réglé les suites à décroissance lente
au milieu des années 1980, Kolountzakis a produit des ensembles de mesure pleine évitant certains
ensembles de Cantor, et des travaux récents étudient des variantes bilipschitziennes et
topologiques, passées en revue dans l'article du cinquantenaire ci-dessous.

\begin{refs}
\item Contexte : Théorème de densité de Lebesgue --- Wikipédia (en anglais) (\url{https://en.wikipedia.org/wiki/Lebesgue%27s_density_theorem})
\item Article : K. J. Falconer, \og{} On a problem of Erdős on sequences and measurable sets \fg{} (1984), Proc. Amer. Math. Soc. 90, 77--78
\item Article : Y. Jung, C.-K. Lai \& Y. Mooroogen, \og{} Fifty years of the Erdős similarity conjecture \fg{} (2024), arXiv:2412.11062 (\url{https://arxiv.org/abs/2412.11062})
\end{refs}

\bigskip

\begin{problem}[{Le problème des ensembles de Furstenberg --- Recherche}]
\textbf{RA-44.} \textbf{Problème.} Fixez $ 0 < s \le 1 $ et $ 0 < t \le 2 $. Un ensemble
$ F \subseteq \mathbb{R}^2 $ est un \textit{ensemble de Furstenberg $ (s,t) $} s'il existe une
famille $ \mathcal{L} $ de droites avec $ \dim_H \mathcal{L} \ge t $, pour la métrique naturelle
sur l'espace des droites, telle que $ \dim_H (F \cap \ell) \ge s $ pour toute
$ \ell \in \mathcal{L} $. La question est de savoir à quel point $ \dim_H F $ peut être
petite.
\begin{enumerate}
\item[(a)] Montrez que le problème contient Kakeya : démontrez qu'un ensemble de Besicovitch du plan
(RA-21) est un ensemble de Furstenberg $ (1,1) $, et vérifiez que la borne conjecturée ci-dessous
redonne alors le théorème de Kakeya plan.
\item[(b)] Démontrez les bornes faciles $ \dim_H F \ge s $ et $ \dim_H F \ge t/2 $ par un argument de
recouvrement, puis lisez et rédigez les deux estimations classiques de Wolff pour les ensembles de
Furstenberg $ (s,1) $, $ \dim_H F \ge 2s $ et $ \dim_H F \ge s + \tfrac12 $, en identifiant
les deux mécanismes différents (un buisson de droites passant par un même point, et un décompte
d'incidences entre tubes) qui les produisent.
\item[(c)] Étudiez le théorème de Ren et Wang (2023), qui règle complètement le cas du plan :
\[ \dim_H F \;\ge\; \min\Big( s+t, \ \frac{3s+t}{2}, \ s+1 \Big), \]
et c'est optimal. Rendez compte de la raison pour laquelle il y a trois termes en concurrence, en
construisant, pour chacun d'eux, un exemple qui le rend contraignant.
\item[(d)] Montez maintenant d'une dimension. Formulez le problème de Furstenberg $ (s,t) $ pour les
$ k $-plans de $ \mathbb{R}^n $, et démontrez une borne optimale dans un cas non couvert par le
théorème plan. \textbf{Statut : ouvert} --- les estimations optimales en dimension supérieure ne sont
pas connues.
\end{enumerate}

\end{problem}

\noindent Le nom vient des travaux de Furstenberg de 1970 sur les intersections
d'ensembles de Cantor et la transversalité ; Katz et Tao ont donné au problème sa forme moderne vers
2001 en montrant qu'il équivaut essentiellement à une conjecture d'anneau discrétisée et qu'il est
étroitement lié au problème de Falconer (RA-23) et à Kakeya (RA-21). Wolff l'avait promu dans les
années 1990 comme cas test propre pour les raisonnements de type Kakeya, et pendant vingt-cinq ans
ses deux bornes ont constitué l'état de l'art. \textbf{Statut en 2026 :} Kevin Ren et Hong Wang ont
entièrement résolu la conjecture plane dans une prépublication de 2023, obtenant au passage une
borne somme-produit de type Elekes et la réponse à une question de projection d'Oberlin ; ce travail
fait partie du cercle d'idées pour lequel Wang a reçu la médaille Fields 2026. Dans
$ \mathbb{R}^n $ pour $ n \ge 3 $, le problème --- sous chacune de ses diverses formulations
naturelles --- reste ouvert.

\begin{refs}
\item Contexte : Ensemble de Besicovitch (ensemble de Kakeya) --- Wikipédia (\url{https://fr.wikipedia.org/wiki/Ensemble_de_Besicovitch})
\item Contexte : Dimension de Hausdorff --- Wikipédia (\url{https://fr.wikipedia.org/wiki/Dimension_de_Hausdorff})
\item Article : K. Ren \& H. Wang, \og{} Furstenberg sets estimate in the plane \fg{} (2023), arXiv:2308.08819 (\url{https://arxiv.org/abs/2308.08819})
\item Lecture : P. Mattila, \textit{Fourier Analysis and Hausdorff Dimension} (Cambridge, 2015)
\end{refs}

\bigskip

\begin{problem}[{Quelles fonctions sont des dérivées ? --- Recherche, short}]
\textbf{RA-45.} \textbf{Problème.} Si $ F : [0,1] \to \mathbb{R} $ est dérivable en tout
point, alors $ f = F' $ est de classe de Baire 1 (limite simple de fonctions continues) et possède
la propriété de Darboux de RA-39 --- mais ces deux conditions réunies ne sont \textbf{pas} suffisantes,
et aucune caractérisation descriptive de la classe des dérivées n'est connue. Construisez une
fonction de classe de Baire 1 vérifiant la propriété de Darboux qui ne soit pas une dérivée, puis
attaquez le problème ouvert lui-même : donnez une liste de conditions intrinsèques sur $ f $, ne
faisant référence à aucun $ F $, qui vaille exactement pour les dérivées.
\end{problem}

\noindent C'est l'une des plus anciennes questions non résolues de la théorie des fonctions
réelles --- Bruckner et Leonard la traitaient déjà comme le problème ouvert central du sujet dans leur
synthèse de 1966, et elle est toujours ouverte soixante ans plus tard. On sait beaucoup de choses sur
ce que les dérivées doivent vérifier : outre la classe de Baire 1 et la propriété de Darboux, il y a
la propriété de Denjoy--Clarkson, selon laquelle $ \{ x : a < f(x) < b \} $ est vide ou
de mesure strictement positive. Mais toute liste proposée de conditions nécessaires s'est jusqu'ici
révélée insuffisante, et la difficulté est instructive : la dérivation est définie par une limite en
chaque point séparément, et personne n'a trouvé l'empreinte globale que ce processus local laisse
derrière lui.

\begin{refs}
\item Contexte : Théorème de Darboux (analyse) --- Wikipédia (\url{https://fr.wikipedia.org/wiki/Th%C3%A9or%C3%A8me_de_Darboux_(analyse)})
\item Contexte : Fonction de Baire --- Wikipédia (\url{https://fr.wikipedia.org/wiki/Fonction_de_Baire})
\item Article : A. M. Bruckner \& J. L. Leonard, \og{} Derivatives \fg{}, \textit{Amer. Math. Monthly} 73 (1966), no 4, partie II, 24--56
\item Lecture : A. M. Bruckner, \textit{Differentiation of Real Functions}
\end{refs}

\bigskip


\section*{Pour aller plus loin}


\vfill
\noindent\emph{Hors de la Caverne} --- des probl\`emes, pas de r\'eponses.
\end{document}
